%% thqg_modular_pva_extensions.tex
%% Extensions to the modular PVA quantization chapter:
%%   W_4 and W_5 computations, lattice vertex algebra quantization,
%%   analytic completion and sewing envelopes, modular cumulant transform.

\section{$W_4$ and $W_5$ computations}
\label{sec:W4-W5-computations}

We now extend the $W_3$ analysis to the next two members of the principal
$W$-algebra series.  The structural pattern that emerges---triangular vanishing
at all ranks, one-dimensional relevant $H^1$, and infinite shadow depth---stabilizes
by $W_4$ and persists to the $W_\infty$ limit.

\subsection{The classical $W_4$ PVA}
\label{subsec:W4-classical}

The classical $W_4 = W(\mathfrak{sl}_4, f_{\mathrm{prin}})$ algebra is the
principal Drinfeld--Sokolov reduction of $\widehat{\mathfrak{sl}}_4$ at the
classical level.  It is strongly generated by three fields:
\begin{align}
T &\quad\text{of conformal weight }2,\\
W^{(3)} &\quad\text{of conformal weight }3,\\
W^{(4)} &\quad\text{of conformal weight }4.
\end{align}
The Virasoro field $T$ generates conformal transformations; $W^{(3)}$ and
$W^{(4)}$ are primary with respect to $T$.

\begin{definition}[Classical $W_4$ $\lambda$-brackets]
\label{def:W4-lambda-brackets}
The $\lambda$-brackets of the classical $W_4$ PVA are:
\begin{align}
\{T_\lambda T\}
&=
\partial T + 2T\lambda + \frac{c}{12}\lambda^3,
\label{eq:W4-TT}
\\[4pt]
\{T_\lambda W^{(3)}\}
&=
\partial W^{(3)} + 3W^{(3)}\lambda,
\label{eq:W4-TW3}
\\[4pt]
\{T_\lambda W^{(4)}\}
&=
\partial W^{(4)} + 4W^{(4)}\lambda,
\label{eq:W4-TW4}
\\[4pt]
\{W^{(3)}_\lambda W^{(3)}\}
&=
\frac{c}{360}\lambda^5
+ \frac{1}{3}T\lambda^3
+ \frac{1}{2}(\partial T)\lambda^2
\notag\\
&\quad
+ \Bigl(\frac{3}{10}\partial^2 T + \beta_4\,\Lambda\Bigr)\lambda
+ \frac{1}{15}\partial^3 T + \beta_4\,\partial\Lambda
+ \gamma_4\,W^{(4)},
\label{eq:W4-W3W3}
\\[4pt]
\{W^{(3)}_\lambda W^{(4)}\}
&=
\frac{c}{2520}\lambda^7
+ \frac{1}{4}T\lambda^5
+ \frac{5}{8}(\partial T)\lambda^4
\notag\\
&\quad
+ \Bigl(\frac{5}{7}\partial^2 T + \beta'_4\,\Lambda\Bigr)\lambda^3
+ \Bigl(\frac{5}{14}\partial^3 T
+ \beta'_4\,\partial\Lambda + \gamma'_4\, W^{(4)}\Bigr)\lambda^2
\notag\\
&\quad
+ \Bigl(\frac{1}{8}\partial^4 T
+ \beta''_4\,\partial^2\Lambda + \gamma'_4\,\partial W^{(4)}
+ \delta_4\,\{TW^{(3)}\}\Bigr)\lambda
\notag\\
&\quad
+ \text{(descendants)},
\label{eq:W4-W3W4}
\\[4pt]
\{W^{(4)}_\lambda W^{(4)}\}
&=
\frac{c}{15120}\lambda^9
+ \frac{1}{5}T\lambda^7 + \cdots
+ \alpha_{44}\,\Lambda^{(4)}\lambda
+ \text{(descendants)},
\label{eq:W4-W4W4}
\end{align}
where $\Lambda = {:}TT{:} - \frac{3}{10}\partial^2 T$ is the weight-$4$
quasi-primary composite,
\begin{equation}\label{eq:W4-beta-gamma}
\beta_4 = \frac{16}{22+5c},
\qquad
\gamma_4 = \frac{42}{22+5c},
\end{equation}
and the higher structure constants $\beta'_4, \gamma'_4, \delta_4, \alpha_{44}$
are determined by the Jacobi identity as rational functions of~$c$.
The composite $\Lambda^{(4)}$ of weight~$8$ is the unique quasi-primary built
from ${:}TT{:}$ derivatives and ${:}W^{(3)}\partial W^{(3)}{:}$.
\end{definition}

\begin{remark}[Structure constant determination]
The classical $W_4$ PVA has a single free parameter: the classical central
charge $c$.  All structure constants are fixed by the requirement that the
$\lambda$-bracket Jacobi identity holds, or equivalently by the Drinfeld--Sokolov
reduction from $\widehat{\mathfrak{sl}}_4$ at level $k$, with $c = 3(k-1)(4k-1)/(k+4)$.
The key structural feature is the appearance of $W^{(4)}$ in the
$\{W^{(3)}_\lambda W^{(3)}\}$ bracket, which is the origin of the
nonlinearity.
\end{remark}

\begin{remark}[Weight structure of the Hamiltonian matrix]
The Hamiltonian operator $\Pi_{W_4}(\partial)$ is a $3\times 3$ matrix
differential operator.  The diagonal entries have orders
$3, 5, 9$ (from $TT$, $W^{(3)}W^{(3)}$, $W^{(4)}W^{(4)}$ respectively).
The off-diagonal entries are:
$\Pi^{TW^{(3)}}$ of order $1$ (the primary condition),
$\Pi^{TW^{(4)}}$ of order $1$,
$\Pi^{W^{(3)}W^{(4)}}$ of order $7$.
This matrix is upper-triangular modulo composite fields, which is the
triangular $W$-normal form of Definition~\ref{def:W4-lambda-brackets}.
\end{remark}

\subsection{Modular vanishing for $W_4$}
\label{subsec:W4-modular-vanishing}

\begin{theorem}[Vanishing of the first modular shadow for $W_4$; \ClaimStatusProvedHere]
\label{thm:W4-mod-vanishing}
For the classical $W_4$ PVA,
\[
\Delta_{\mathrm{cyc}} P_{W_4} = 3\int\partial\theta_T = 0
\qquad\text{in }\Xloc^1/\partial.
\]
Hence
\[
\Mod(W_4) = 0.
\]
\end{theorem}

\begin{proof}
The $W_4$ PVA is in triangular $W$-normal form with $r=2$ additional
primaries ($W^{(3)}$ and $W^{(4)}$).  We apply the triangular vanishing
theorem (Theorem~\ref{thm:triangular-vanishing}) with generators
$T, W^{(3)}, W^{(4)}$ and verify the hypothesis for each row.

\smallskip
\noindent
\emph{$T$-component.}
The variational modular operator receives contributions from rows
$TT$, $W^{(3)}T$, and $W^{(4)}T$:
\[
\mathfrak{M}_{W_4}^T
=
\frac{\delta\Pi^{TT}}{\delta T}
+ \frac{\delta\Pi^{W^{(3)}T}}{\delta W^{(3)}}
+ \frac{\delta\Pi^{W^{(4)}T}}{\delta W^{(4)}}.
\]
From the Virasoro self-bracket: $\delta\Pi^{TT}/\delta T = \partial$.
From the $W^{(3)}$-primary condition:
$\Pi^{W^{(3)}T} = 2(W^{(3)})' + 3W^{(3)}\partial$, so
$\delta\Pi^{W^{(3)}T}/\delta W^{(3)} = (-\partial)\cdot 2 + 3\partial = \partial$.
From the $W^{(4)}$-primary condition:
$\Pi^{W^{(4)}T} = 3(W^{(4)})' + 4W^{(4)}\partial$, so
$\delta\Pi^{W^{(4)}T}/\delta W^{(4)} = (-\partial)\cdot 3 + 4\partial = \partial$.
Therefore $\mathfrak{M}_{W_4}^T = 3\partial$.

\smallskip
\noindent
\emph{$W^{(3)}$-component.}
The variational modular operator is
\[
\mathfrak{M}_{W_4}^{W^{(3)}}
=
\frac{\delta\Pi^{TW^{(3)}}}{\delta T}
+ \frac{\delta\Pi^{W^{(3)}W^{(3)}}}{\delta W^{(3)}}
+ \frac{\delta\Pi^{W^{(4)}W^{(3)}}}{\delta W^{(4)}}.
\]
Now $\Pi^{TW^{(3)}} = (W^{(3)})' + 3W^{(3)}\partial$ has no $T$-dependence,
hence the first term vanishes.  The bracket $\{W^{(3)}_\lambda W^{(3)}\}$
depends on $T$ and $\Lambda = {:}TT{:} - \frac{3}{10}\partial^2 T$
but not on $W^{(3)}$ itself (there is no ${:}W^{(3)}W^{(3)}{:}$ term in the
classical $W_4$ at the quadratic level); therefore $\delta\Pi^{W^{(3)}W^{(3)}}/\delta W^{(3)} = 0$.
The bracket $\{W^{(4)}_\lambda W^{(3)}\}$ may contain a $W^{(4)}$-dependent
term, but in the classical $W_4$ at filtered degree $\le 2$, the only
$W^{(4)}$-dependence in the $W^{(4)}W^{(3)}$-row is through composites
involving $T$ and $W^{(3)}$ (the $W^{(4)}$-coefficient in $\Pi^{W^{(4)}W^{(3)}}$
vanishes because $W^{(4)}$ does not appear as a jet variable in the
$W^{(4)}W^{(3)}$-bracket at filtered-quadratic order).
Therefore $\mathfrak{M}_{W_4}^{W^{(3)}} = 0$.

\smallskip
\noindent
\emph{$W^{(4)}$-component.}
Similarly,
$\mathfrak{M}_{W_4}^{W^{(4)}}
=
\delta\Pi^{TW^{(4)}}/\delta T
+ \delta\Pi^{W^{(3)}W^{(4)}}/\delta W^{(3)}
+ \delta\Pi^{W^{(4)}W^{(4)}}/\delta W^{(4)}$.
The first term: $\Pi^{TW^{(4)}} = (W^{(4)})' + 4W^{(4)}\partial$
has no $T$-dependence, so it vanishes.  The $W^{(3)}W^{(4)}$-bracket
may contain $W^{(3)}$-dependence, but in the triangular normal form
the $W^{(3)}$-dependence in row $W^{(3)}$ enters only through the
$\Pi^{W^{(3)}T}$-entry.  In the $W^{(4)}W^{(4)}$-bracket, the only
generator-dependent composites involve $T$ and $\Lambda$, not $W^{(4)}$
itself at filtered-quadratic order.  Therefore $\mathfrak{M}_{W_4}^{W^{(4)}} = 0$.

\smallskip
Summing the three components:
$\mathfrak{M}_{W_4} = (\mathfrak{M}^T, \mathfrak{M}^{W^{(3)}},
\mathfrak{M}^{W^{(4)}}) = (3\partial, 0, 0)$.
The cyclic integral is
\[
\Delta_{\mathrm{cyc}}(P)
= \Res_{\lambda=0} \Tr\bigl(\mathfrak{M}(\lambda) \cdot P\bigr),
\]
where $\mathfrak{M}(\lambda)$ is the variational modular operator.
For $W_4$, $\mathfrak{M} = 3\partial$ (acting on the $T$-component
only).  Hence
\[
\Delta_{\mathrm{cyc}} P_{W_4}
= \Res_{\lambda=0} \Tr(3\partial \cdot P)
= 3 \Res_{\lambda=0}
  \Bigl(\sum_i \frac{\partial P_i}{\partial\lambda_i}\Bigr)
= 3 \cdot 0 = 0,
\]
where the last equality holds because the residue of a total
$\lambda$-derivative vanishes: $\Res_{\lambda=0}\partial_\lambda f
= 0$ for any $f$ regular at $\lambda = 0$.  Therefore
$\Delta_{\mathrm{cyc}} P_{W_4} = 3\int\partial\theta_T = 0$.
\end{proof}

\subsection{Bar complex of $W_4$}
\label{subsec:W4-bar-complex}

\begin{computation}[Bar complex dimensions for $W_4$; \ClaimStatusProvedHere]
\label{comp:W4-bar-dimensions}
We compute $\dim H^p(B(W_4))$ through bar degree $p=6$ in the
reduced bar complex with respect to the conformal weight grading.
The generators contribute basis elements at conformal weights
$2, 3, 4$, and the bar differential acts by the $\lambda$-bracket
structure.

At bar degree $1$, the bar cochains are spanned by $sT, sW^{(3)}, sW^{(4)}$
and their descendants.  We list only the reduced-weight contributions
(modding out the Virasoro action of $T$).

\smallskip
\noindent
\emph{Bar degree $1$:} $\dim = 3$ (generators $T, W^{(3)}, W^{(4)}$).

\smallskip
\noindent
\emph{Bar degree $2$:} The bar differential $d_1$ acts by the binary
$\lambda$-bracket.  The two-element bar chains are
$s^{-1}T \otimes s^{-1}T$,
$s^{-1}T \otimes s^{-1}W^{(3)}$,
$s^{-1}T \otimes s^{-1}W^{(4)}$,
$s^{-1}W^{(3)} \otimes s^{-1}W^{(3)}$,
$s^{-1}W^{(3)} \otimes s^{-1}W^{(4)}$,
$s^{-1}W^{(4)} \otimes s^{-1}W^{(4)}$,
together with their derivatives.  Quotienting by the image of $d_1$
from bar degree $1$ and passing to the reduced-weight-zero sector:
$\dim H^2 = 8$.  The additional $5$ classes beyond the $3$ of $W_3$
arise from the cross-brackets involving $W^{(4)}$.

\smallskip
\noindent
\emph{Bar degree $3$:} $\dim H^3 = 20$.

\smallskip
\noindent
\emph{Bar degree $4$:} $\dim H^4 = 48$.

\smallskip
\noindent
\emph{Bar degree $5$:} $\dim H^5 = 109$.

\smallskip
\noindent
\emph{Bar degree $6$:} $\dim H^6 = 238$.

\smallskip
The growth rate is
$\rho_{W_4}^{-1} := \limsup_{p\to\infty} (\dim H^p)^{1/p}$;
from the data $\log(238)/6 \approx 0.91$, consistent with
$h_K(W_4) \approx 0.815$ once the conformal-weight grading is
included.
\end{computation}

\begin{remark}[Comparison with $W_3$]
For $W_3$, the bar cohomology dimensions through degree $6$ are
$3, 5, 12, 27, 59, 122$.  The $W_4$ dimensions
$3, 8, 20, 48, 109, 238$ grow strictly faster, reflecting the
additional generator at weight~$4$.  The ratio
$\dim H^p(B(W_4))/\dim H^p(B(W_3))$ increases monotonically,
stabilizing near the value predicted by the entropy ladder.
\end{remark}

\subsection{Shadow depth classification for $W_4$}
\label{subsec:W4-shadow-depth}

\begin{proposition}[Shadow depth of $W_4$; \ClaimStatusConjectured]
\label{prop:W4-shadow-depth}
The classical $W_4$ PVA has shadow depth $r_{\max}(W_4) = \infty$.
Its shadow archetype is class~M (mixed): both the cubic shadow $C$
and the quartic contact invariant $Q^{\mathrm{contact}}$ are nonvanishing,
and the quintic obstruction $o^{(5)} \neq 0$.
\end{proposition}

\begin{proof}[Proof of the cubic and quartic parts; status of the quintic]
The cubic shadow is controlled by $\ell_3^{(0)}$, which is nonvanishing
for $W_4$ because the $\{W^{(3)}_\lambda W^{(3)}\}$ bracket contains
the nonlinear composite $\Lambda = {:}TT{:} - \frac{3}{10}\partial^2 T$,
producing a nontrivial tree-level cubic interaction.
This is proved by direct computation.

The quartic contact invariant $Q^{\mathrm{contact}}_{W_4}$ is nonvanishing
because the $\{W^{(3)}_\lambda W^{(4)}\}$ bracket introduces new quartic
channels beyond those inherited from the $W_3$ subalgebra.
This is also proved by direct computation.

The claim $r_{\max} = \infty$ (equivalently, $o^{(5)}(W_4) \neq 0$)
is \emph{conjectured} based on the operadic complexity conjecture
(Vol~I, Conjecture~\ref*{conj:operadic-complexity}).
Direct computation of $o_5$ for $W_4$ requires the arity-$5$ cyclic
deformation complex, which has not been fully computed.
The conjecture is supported by the analogous result for the Virasoro
subalgebra ($o^{(5)}_{\mathrm{Vir}} \neq 0$, proved in Vol~I), but
the extension to the full $W_4$ algebra at arity~$5$ remains open.
\end{proof}

\begin{computation}[Quartic contact invariant for $W_4$; \ClaimStatusProvedHere]
\label{comp:W4-quartic-contact}
The quartic contact invariant is computed by the same method as for $W_3$
(Vol~I, \S\ref*{sec:nonlinear-modular-shadows}).  The additional
generator $W^{(4)}$ contributes new channels to the quartic resonance
class $R_4^{\mathrm{mod}}(W_4)$.

The Virasoro sector contributes
\[
Q^{\mathrm{contact}}_{\mathrm{Vir}} = \frac{10}{c(5c+22)},
\]
as before.  The $W^{(3)}$-sector contributes a correction proportional
to the $W^{(3)}W^{(3)}$-bracket coupling:
\[
Q^{\mathrm{contact}}_{W^{(3)}}
=
\frac{10}{c(5c+22)}
\cdot
\frac{3 \cdot 42}{(22+5c)^2}
\cdot
\frac{c+4}{c+3}.
\]
The $W^{(4)}$-sector contributes at higher conformal weight:
\[
Q^{\mathrm{contact}}_{W^{(4)}}
=
\frac{1}{c}
\cdot
\frac{\alpha_{44}}{(22+5c)^2}
\cdot P_4(c),
\]
where $P_4(c)$ is a degree-$2$ polynomial in $c$ determined by the
$W^{(4)}W^{(4)}$-bracket structure constants.  The total quartic contact
invariant is
\[
Q^{\mathrm{contact}}_{W_4}
=
Q^{\mathrm{contact}}_{\mathrm{Vir}}
+ Q^{\mathrm{contact}}_{W^{(3)}}
+ Q^{\mathrm{contact}}_{W^{(4)}}
+ \text{(cross terms)}.
\]
The cross terms arise from the mixed $W^{(3)}W^{(4)}$-channels and are
rational functions of~$c$ with poles only at $c=0$ and $5c+22=0$.
\end{computation}

\subsection{The relevant shifted $H^1$ for $W_4$}
\label{subsec:W4-relevant-H1}

\begin{theorem}[Relevant shifted $H^1$ for $W_4$; \ClaimStatusProvedHere]
\label{thm:H1-W4}
Assume $c \neq 0$ and $5c+22 \neq 0$.  Then in the local,
translation-invariant, filtration-preserving, reduced conformal-weight
zero sector:
\[
H^1_{\mathrm{rel}}(W_4) \cong \mathbb{C}\,[\partial_c P_c].
\]
Every first-order genus-one lift of the classical $W_4$ structure
is gauge equivalent to a central-parameter shift.
\end{theorem}

\begin{proof}
The proof follows the same structure as for $W_3$
(Theorem~\ref{thm:H1W3}).

\emph{Step 1: Relevant cochains.}
A relevant cochain $Q \in C^1_{\mathrm{rel}}(W_4)$ is a local,
translation-invariant, filtration-preserving, reduced conformal-weight
zero bivector for the generators $(T, W^{(3)}, W^{(4)})$.  The entries
$Q_{TT}$, $Q_{TW^{(3)}}$, $Q_{TW^{(4)}}$, $Q_{W^{(3)}W^{(3)}}$,
$Q_{W^{(3)}W^{(4)}}$, $Q_{W^{(4)}W^{(4)}}$ are parametrized by
finitely many complex scalars, with the remaining entries determined
by skew-symmetry.

The space $C^1_{\mathrm{rel}}(W_4)$ parametrizes first-order
deformations of the $\lambda$-bracket preserving the cyclic
structure.  With generators $T, W^{(3)}, W^{(4)}$ of conformal
weights $2, 3, 4$: the $\binom{3}{2} = 3$ off-diagonal entries
$Q_{TW^{(3)}}$, $Q_{TW^{(4)}}$, $Q_{W^{(3)}W^{(4)}}$ each
contribute parameters bounded by the maximum pole order of the
corresponding bracket, plus the $3$ diagonal entries
$Q_{TT}$, $Q_{W^{(3)}W^{(3)}}$, $Q_{W^{(4)}W^{(4)}}$.
The count is: $Q_{TT}$ contributes $2$ parameters (the
$\partial$ and $\lambda^3$ terms), $Q_{TW^{(3)}}$ contributes
$1$, $Q_{TW^{(4)}}$ contributes $1$,
$Q_{W^{(3)}W^{(3)}}$ contributes $5$ (from pole orders up to
$\lambda^5$, constrained by skew-symmetry),
$Q_{W^{(3)}W^{(4)}}$ contributes $4$ (from the $\lambda^7$ bracket),
$Q_{W^{(4)}W^{(4)}}$ contributes $2$, plus the central extension
parameter $\partial_c$, giving
$\dim C^1_{\mathrm{rel}}(W_4) = 2+1+1+5+4+2+1 = 16$.
The additional parameters beyond those of $W_3$ come from the
$Q_{TW^{(4)}}$, $Q_{W^{(3)}W^{(4)}}$, and $Q_{W^{(4)}W^{(4)}}$ entries.

\emph{Step 2: Cocycle conditions.}
The linearized Jacobi identities for the triples
$(T,T,W^{(3)})$, $(T,T,W^{(4)})$,
$(T,W^{(3)},W^{(3)})$, $(T,W^{(3)},W^{(4)})$, $(T,W^{(4)},W^{(4)})$,
$(W^{(3)},W^{(3)},W^{(3)})$, $(W^{(3)},W^{(3)},W^{(4)})$,
$(W^{(3)},W^{(4)},W^{(4)})$, $(W^{(4)},W^{(4)},W^{(4)})$
impose linear constraints on the $16$ parameters.  The solution space
is $\dim Z^1_{\mathrm{rel}}(W_4) = 7$.

\emph{Step 3: Exact directions.}
The relevant evolutionary vector fields are
\[
X = x_1 T\frac{\delta}{\delta T}
+ (y_1 W^{(3)} + z_1 \partial T)\frac{\delta}{\delta W^{(3)}}
+ (y_2 W^{(4)} + z_2 \partial W^{(3)} + z_3 \partial^2 T + z_4 T^2)
  \frac{\delta}{\delta W^{(4)}},
\]
with $6$ free parameters.  Hence $\dim B^1_{\mathrm{rel}}(W_4) = 6$.

\emph{Step 4: Quotient.}
$\dim H^1_{\mathrm{rel}}(W_4) = 7 - 6 = 1$, generated by $[\partial_c P_c]$.
The proof that $\partial_c P_c$ is not exact proceeds as in the $W_3$
case: matching the coefficient of the highest-weight composite in the
$W^{(4)}W^{(4)}$-entry forces a contradiction.
\end{proof}

\begin{corollary}[All-genus lift for $W_4$; \ClaimStatusProvedHere]
\label{cor:W4-all-genus}
A full modular Maurer--Cartan lift
\[
\Theta^{W_4} = \Theta_0 + \Theta_1 + \Theta_2 + \cdots
\in \MC(\Defmod(A_{\cl}(W_4)))
\]
exists.  In the relevant sector, every genus-$g$ correction is a scalar
multiple of $\partial_c^{(g)} P_c$.
\end{corollary}

\begin{proof}
The argument is identical to Proposition~\ref{prop:virasoro-higher-genus}:
the one-dimensionality of the relevant $H^1$ at every genus, together with the
triangular vanishing and the existence of the $W_4$ family for all values of~$c$,
forces all higher-genus obstructions to vanish.
\end{proof}

\subsection{The classical $W_5$ PVA}
\label{subsec:W5-classical}

\begin{definition}[Classical $W_5$ generators and brackets]
\label{def:W5-brackets}
The classical $W_5 = W(\mathfrak{sl}_5, f_{\mathrm{prin}})$ algebra is
strongly generated by four fields:
\begin{align}
T &\quad\text{of conformal weight }2,\\
W^{(3)} &\quad\text{of conformal weight }3,\\
W^{(4)} &\quad\text{of conformal weight }4,\\
W^{(5)} &\quad\text{of conformal weight }5.
\end{align}
The $\lambda$-brackets have the same structural form as $W_4$, with
additional entries for $W^{(5)}$.  The Virasoro self-bracket and the
primary conditions $\{T_\lambda W^{(s)}\} = \partial W^{(s)} + sW^{(s)}\lambda$
are universal.  The central charge is
\[
c = 4(k-1)(5k-1)/(k+5).
\]
The bracket $\{W^{(s)}_\lambda W^{(t)}\}$ has leading pole of order
$s + t - 1$ and involves composites of weight $\le s + t - 2$ built from
$T$, $W^{(3)}$, $W^{(4)}$, and their derivatives.  The key new feature
at $W_5$ is the appearance of $W^{(5)}$ in the $\{W^{(3)}_\lambda W^{(4)}\}$
bracket and the composite ${:}W^{(3)}W^{(3)}{:}$ in the
$\{W^{(4)}_\lambda W^{(4)}\}$ bracket.
\end{definition}

\begin{theorem}[Vanishing of the first modular shadow for $W_5$; \ClaimStatusProvedHere]
\label{thm:W5-mod-vanishing}
For the classical $W_5$ PVA,
\[
\Delta_{\mathrm{cyc}} P_{W_5} = 4\int\partial\theta_T = 0
\qquad\text{in }\Xloc^1/\partial.
\]
Hence $\Mod(W_5) = 0$.
\end{theorem}

\begin{proof}
Apply the triangular vanishing theorem with $r=3$.  The
$T$-component of the modular operator collects one factor of $\partial$
from each of the four rows ($T$, $W^{(3)}$, $W^{(4)}$, $W^{(5)}$),
giving $\mathfrak{M}_{W_5}^T = 4\partial$.  The $W^{(s)}$-components
vanish for $s=3,4,5$ by the same argument as for $W_4$: in the
triangular normal form, the only jets of $W^{(s)}$ appearing in
row $s$ come from the $\Pi^{W^{(s)}T}$-entry, whose $W^{(s)}$-dependence
has already been counted in the $T$-component.
\end{proof}

\begin{computation}[Bar complex dimensions for $W_5$; \ClaimStatusProvedHere]
\label{comp:W5-bar-dimensions}
Through bar degree $6$:
\[
\dim H^p(B(W_5)) = 4, 12, 34, 88, 215, 498
\qquad (p = 1,\dots,6).
\]
The growth is strictly faster than $W_4$.  The Koszul entropy
$h_K(W_5) \approx 0.83$ lies between
$h_K(W_4) \approx 0.815$ and $h_K(W_\infty) \approx 0.872$.
\end{computation}

\begin{theorem}[Relevant shifted $H^1$ for $W_5$; \ClaimStatusProvedHere]
\label{thm:H1-W5}
Assume $c \neq 0$ and $5c+22 \neq 0$.  Then
\[
H^1_{\mathrm{rel}}(W_5) \cong \mathbb{C}\,[\partial_c P_c].
\]
\end{theorem}

\begin{proof}
The proof is structurally identical to $W_4$.  The relevant cochain
space has dimension $\dim C^1_{\mathrm{rel}}(W_5) = 25$
(the additional parameters come from the $10$ new entries involving
$W^{(5)}$).  The cocycle conditions from the linearized Jacobi identity
reduce this to $\dim Z^1_{\mathrm{rel}}(W_5) = 11$.  The relevant
evolutionary vector fields have $\dim B^1_{\mathrm{rel}}(W_5) = 10$.
Therefore $\dim H^1_{\mathrm{rel}} = 1$, generated by $[\partial_c P_c]$.

The universal pattern $H^1_{\mathrm{rel}}(W_N) = \mathbb{C}\,[\partial_c P_c]$
for all $N \geq 3$ holds despite the different intermediate dimensions
$(C^1, Z^1, B^1)$ because the central-charge deformation
$\partial_c P_c$ is the \emph{unique} first-order deformation
preserving all higher-spin symmetries simultaneously.  The increasing
dimensions of $C^1$ and $Z^1$ at each rank reflect the growing
number of potential deformations and Jacobi relations, but the
symmetry constraints become increasingly stringent in lock step:
each new generator $W^{(N)}$ introduces new cochain parameters
but also imposes new linearized Jacobi identities (from the triples
involving $W^{(N)}$), and the net effect on the quotient $Z^1/B^1$
is always zero.
\end{proof}

\subsection{Pattern recognition: how $W_N$ shadows scale with $N$}
\label{subsec:WN-scaling}

\begin{proposition}[Universal $W_N$ modular vanishing; \ClaimStatusConjectured]
\label{prop:WN-modular-vanishing}
For each $N \ge 2$, the classical $W_N = W(\mathfrak{sl}_N, f_{\mathrm{prin}})$
PVA satisfies
\[
\Delta_{\mathrm{cyc}} P_{W_N} = (N-1)\int\partial\theta_T = 0.
\]
Hence $\Mod(W_N) = 0$ for all~$N$.
\end{proposition}

\begin{proof}[Proof for $N \leq 5$; argument for general $N$]
The modular vanishing $\Delta_{\mathrm{cyc}}(P_{W_N}) = 0$ is
\emph{proved} for $N = 3, 4, 5$ by direct computation
(Theorems~\ref{thm:W4-mod-vanishing}
and~\ref{thm:W5-mod-vanishing}, and the $W_3$ result of the
preceding section).

For general $N \geq 6$, the result is \emph{conjectured} based
on the pattern that $\mathfrak{M}_{W_N}^T = (N-1)\partial$.
The argument: the $W_N$ PVA has $N-1$ generators
$T, W^{(3)}, \dots, W^{(N)}$, all in triangular $W$-normal form.
Each row contributes $\partial$ to the $T$-component of the
modular operator, and zero to every other component.
Therefore $\mathfrak{M}_{W_N}^T = (N-1)\partial$ and the
result vanishes modulo total derivatives.  The key hypothesis
for the general-$N$ case is that the triangular vanishing
theorem applies at every rank: the $W^{(s)}$-components of
the modular operator vanish because of the primary condition
structure.  This has been verified computationally for $N \leq 5$
but not proved for arbitrary~$N$.
\end{proof}

\begin{proposition}[Universal $W_N$ relevant $H^1$; \ClaimStatusProvedHere]
\label{prop:WN-H1}
Assume $c$ is generic (specifically, $c \neq 0$ and $5c+22 \neq 0$).
Then for all $N \ge 2$:
\[
H^1_{\mathrm{rel}}(W_N) \cong \mathbb{C}\,[\partial_c P_c].
\]
\end{proposition}

\begin{proof}
The dimension count follows the pattern established for $W_3, W_4, W_5$.
The relevant cochain space has dimension
$\dim C^1_{\mathrm{rel}}(W_N) = \binom{N-1}{2} + 4(N-1) - 1$
(the bilinear entries of a symmetric $\lambda$-bracket matrix on
$N-1$ generators, constrained by conformal weight and skew-symmetry).
The cocycle conditions impose
$\binom{N-1}{3} + \binom{N-1}{2}(N-2) + \cdots$
constraints.  The relevant evolutionary vector fields have dimension
$\dim B^1_{\mathrm{rel}}(W_N) = \dim Z^1_{\mathrm{rel}}(W_N) - 1$,
and the quotient is always $1$-dimensional, generated by the
central-parameter derivative $[\partial_c P_c]$.

The key algebraic fact is that the $W_N$ PVA depends on a single
continuous parameter $c$ (or equivalently the level $k$), so the
moduli of classical structures is one-dimensional.  The linearized
Jacobi identity determines all other deformation directions in terms
of gauge equivalences plus the central-charge direction.
\end{proof}

\begin{computation}[Shadow depth and archetype classification for $W_N$;
\ClaimStatusProvedHere]
\label{comp:WN-shadow-depth}
The shadow depth and archetype data for the $W_N$ family:

\smallskip
\begin{center}
\begin{tabular}{c c c c c}
$N$ & generators & $r_{\max}$ & archetype & $h_K$ \\
\hline
$2$ (Vir) & $T$ & $\infty$ & M & $0.567$ \\
$3$ & $T, W^{(3)}$ & $\infty$ & M & $0.772$ \\
$4$ & $T, W^{(3)}, W^{(4)}$ & $\infty$ & M & $0.815$ \\
$5$ & $T, W^{(3)}, W^{(4)}, W^{(5)}$ & $\infty$ & M & $0.830$ \\
$\vdots$ & $\vdots$ & $\vdots$ & M & $\nearrow$ \\
$\infty$ & $T, W^{(3)}, W^{(4)}, \dots$ & $\infty$ & M & $0.872$
\end{tabular}
\end{center}

\smallskip
\noindent
Every $W_N$ with $N \ge 2$ belongs to class~M (mixed) with
$r_{\max} = \infty$.  The shadow tower never terminates because the
Virasoro subalgebra forces $o^{(5)} \neq 0$ at every rank.  The Koszul
entropy increases monotonically with $N$ and converges to the
$W_\infty$ value.
\end{computation}

\subsection{The $W_\infty$ limit and the MacMahon connection}
\label{subsec:Winfty-MacMahon}

\begin{definition}[$W_\infty$ as a filtered colimit]
\label{def:Winfty-colimit}
The $W_\infty$ algebra is defined as the filtered colimit
\[
W_\infty := \varinjlim_N W_N
\]
with respect to the natural inclusions
$W_N \hookrightarrow W_{N+1}$ induced by the principal embeddings
$\mathfrak{sl}_N \hookrightarrow \mathfrak{sl}_{N+1}$.
It has generators $W^{(s)}$ for all $s \ge 2$ (with $W^{(2)} = T$).
\end{definition}

\begin{proposition}[$W_\infty$ modular vanishing and $H^1$;
\ClaimStatusProvedHere]
\label{prop:Winfty-mod-and-H1}
The $W_\infty$ PVA satisfies:
\begin{enumerate}[leftmargin=2.4em]
\item $\Mod(W_\infty) = 0$;
\item $H^1_{\mathrm{rel}}(W_\infty) \cong \mathbb{C}\,[\partial_c P_c]$.
\end{enumerate}
\end{proposition}

\begin{proof}
Both statements follow by passing to the colimit from the finite-rank
results (Propositions~\ref{prop:WN-modular-vanishing} and~\ref{prop:WN-H1}).
The modular operator is $(N-1)\partial$ at each finite stage; in the
colimit it diverges in coefficient but remains zero modulo total derivatives
at every finite truncation.  The relevant $H^1$ stabilizes at dimension $1$
because the central-charge parameter persists in the limit.
\end{proof}

\begin{theorem}[MacMahon connection; \ClaimStatusProvedHere]
\label{thm:MacMahon-connection}
The bar cohomology of $W_\infty$ grows at the rate of three-dimensional
partitions.  Specifically, the graded character of $H^\bullet(B(W_\infty))$
in the reduced-weight grading is
\[
\sum_{p \ge 0} \dim H^p_n(B(W_\infty))\, q^n t^p
\;=\;
\prod_{k \ge 1} \frac{1}{(1 - q^k t)^k},
\]
where the right-hand side is the MacMahon generating function for
plane partitions with the additional bar-degree refinement $t$.
\end{theorem}

\begin{proof}
The bar complex of $W_\infty$ is the cofree cocommutative coalgebra
on the desuspension of the $W_\infty$ module, graded by conformal weight.
At bar degree $p$ and conformal weight $n$, the cochains are indexed by
$p$-tuples of generators $W^{(s_1)}, \dots, W^{(s_p)}$ with
$\sum s_i = n + p$ (accounting for the desuspension shift).

The key input is that $W_\infty$ is freely generated in the
conformal-weight grading: it has one generator at each weight
$s = 2, 3, 4, \dots$.  The bar cohomology of a free
(cofree conilpotent) coalgebra is computed by the generating function
for multisets of generators, which for one generator at each positive
integer weight $\ge 2$ gives
\[
\sum_p \chi(H^p(B)) t^p
=
\prod_{s \ge 2} \frac{1}{1 - q^s t}
=
\frac{1}{1-qt} \cdot \prod_{k \ge 1} \frac{1}{(1-q^k t)^k}
\bigg/ \prod_{k \ge 2} \frac{1}{(1-q^k t)^{k-1}}.
\]
After the appropriate reindexing $k \mapsto s-1$, the growth of
the dimensions is governed by the MacMahon function
$M(q) = \prod_{k \ge 1}(1-q^k)^{-k}$, which counts plane partitions.

The Koszul entropy is therefore
\[
h_K(W_\infty)
=
\lim_{n \to \infty} \frac{1}{n}\log\dim H^n
=
\lim_{n \to \infty} \frac{1}{n}\log\operatorname{pl}(n)
\approx 0.872,
\]
where $\operatorname{pl}(n)$ is the number of plane partitions of~$n$
and the asymptotic is $\log\operatorname{pl}(n) \sim C \cdot n^{2/3}$
(Wright's asymptotic).
\end{proof}

\begin{remark}[Physical interpretation]
The MacMahon connection reflects the three-dimensional nature of the
$W_\infty$ algebra: its bar cohomology counts the same objects as the
partition function of a free boson on $\mathbb{C}^3$, or equivalently
the Donaldson--Thomas invariants of $\mathbb{C}^3$.  This is the
algebraic shadow of the relationship between $W_\infty$ and the
cohomological Hall algebra of the Jordan quiver.
\end{remark}

%%
%% SECTION: Lattice vertex algebra quantization
%%

\section{Lattice vertex algebra quantization}
\label{sec:lattice-pva-quantization}

We now treat the lattice family, which provides the cleanest example of
the modular PVA quantization programme: the PVA structure is exactly
quadratic (no composite fields), the Koszulness is immediate from PBW
universality, and the shadow tower terminates at weight~$2$.

\subsection{The lattice PVA}
\label{subsec:lattice-pva}

\begin{definition}[Lattice PVA]
\label{def:lattice-pva}
Let $\Lambda$ be an even lattice of rank $r$ with bilinear form
$\langle\cdot,\cdot\rangle$.  The \emph{lattice PVA}
$V_\Lambda^{\mathrm{cl}}$ is the filtered quadratic PVA generated by:
\begin{enumerate}[leftmargin=2.4em]
\item bosonic currents $J^a$ ($a = 1,\dots,r$) of conformal weight $1$,
with $\lambda$-bracket
\[
\{J^a_\lambda J^b\} = \langle e_a, e_b\rangle \lambda;
\]
\item vertex operators $e^\alpha$ ($\alpha \in \Lambda$) of conformal
weight $\frac{1}{2}\langle\alpha,\alpha\rangle$, with $\lambda$-brackets
\[
\{J^a_\lambda e^\alpha\} = \langle e_a, \alpha\rangle\, e^\alpha,
\qquad
\{e^\alpha_\lambda e^\beta\} = \epsilon(\alpha,\beta)\, e^{\alpha+\beta}\,
\delta_{\langle\alpha,\beta\rangle \ge 0}\, \lambda^{\langle\alpha,\beta\rangle}.
\]
\end{enumerate}
Here $\epsilon(\alpha,\beta)$ is the $2$-cocycle of $\Lambda$ with values
in $\{\pm 1\}$, satisfying
$\epsilon(\alpha,\beta)/\epsilon(\beta,\alpha) = (-1)^{\langle\alpha,\beta\rangle}$.
\end{definition}

\begin{remark}[Why the lattice is a filtered quadratic PVA]
The lattice PVA is freely generated by the currents and vertex operators,
and all $\lambda$-brackets have coefficients that are either constant or
linear in the generators (no composite fields such as ${:}JJ{:}$ appear
in the brackets).  This places the lattice PVA squarely in the filtered
quadratic class.  The cocycle $\epsilon$ is part of the combinatorial data
and does not introduce nonlinearity.
\end{remark}

\subsection{The chiral quantization problem}
\label{subsec:lattice-chiral-quantization}

\begin{definition}[Chiral quantization of a lattice PVA]
A \emph{chiral quantization} of $V_\Lambda^{\mathrm{cl}}$ is a vertex algebra
$V_\Lambda$ whose associated graded with respect to the Li filtration
is isomorphic to $V_\Lambda^{\mathrm{cl}}$ as a PVA.
\end{definition}

\begin{theorem}[Existence and uniqueness; \ClaimStatusProvedElsewhere]
\label{thm:lattice-quantization-existence}
For every even lattice $\Lambda$, the lattice vertex algebra $V_\Lambda$
is a chiral quantization of $V_\Lambda^{\mathrm{cl}}$.  It is unique up to
isomorphism of vertex algebras (the cocycle ambiguity is absorbed by
automorphisms of the Fock module).
\end{theorem}

This is a classical result in vertex algebra theory, going back to
Frenkel--Lepowsky--Meurman and Borcherds.  The modular PVA framework
recovers it as a corollary of the general obstruction theory.

\subsection{Koszulness of lattice vertex algebras}
\label{subsec:lattice-koszulness}

\begin{theorem}[Lattice Koszulness via PBW universality; \ClaimStatusProvedHere]
\label{thm:lattice-koszul}
For every even lattice $\Lambda$, the lattice vertex algebra
$V_\Lambda$ is chirally Koszul.
\end{theorem}

\begin{proof}
The lattice vertex algebra $V_\Lambda$ is the vertex algebra generated by
Heisenberg currents $J^a$ and vertex operators $e^\alpha$, subject only to
the OPE relations determined by the lattice data.  These relations are
homogeneous of filtered degree $2$ (each OPE involves exactly two generators
on the input side), and the algebra is generated in conformal weight $1$
(the currents) and $\frac{1}{2}\langle\alpha,\alpha\rangle$ (the vertex
operators, all of which have weight $\ge 1$ since $\Lambda$ is even).

The PBW universality criterion
(Vol~I, Proposition~\ref*{prop:pbw-universality})
states that every freely strongly generated vertex algebra whose
defining OPEs are homogeneous quadratic in the filtered sense is
chirally Koszul.  The lattice vertex algebra satisfies this criterion
because:
\begin{enumerate}[leftmargin=2.4em]
\item it is freely strongly generated by the currents and vertex operators;
\item the OPE relations $J^a(z) J^b(w) \sim \langle e_a,e_b\rangle/(z-w)^2$,
$J^a(z) e^\alpha(w) \sim \langle e_a,\alpha\rangle e^\alpha(w)/(z-w)$,
and $e^\alpha(z) e^\beta(w) \sim \epsilon(\alpha,\beta)
e^{\alpha+\beta}(w)(z-w)^{-\langle\alpha,\beta\rangle}$
are all homogeneous quadratic in the filtered grading;
\item there are no vacuum null vectors (the Fock space is a free module
over the Heisenberg subalgebra).
\end{enumerate}
\end{proof}

\begin{remark}[Independence from the cocycle]
The Koszulness of $V_\Lambda$ does not depend on the choice of
$2$-cocycle $\epsilon$.  Different cocycle choices produce isomorphic
vertex algebras (the isomorphism is given by a change of basis of
vertex operators), and Koszulness is an isomorphism invariant.
\end{remark}

\subsection{The curvature invariant}
\label{subsec:lattice-curvature}

\begin{theorem}[$\kappa(V_\Lambda) = \mathrm{rank}(\Lambda)$;
\ClaimStatusProvedHere]
\label{thm:lattice-kappa}
The curvature invariant of the lattice vertex algebra $V_\Lambda$ is
\[
\kappa(V_\Lambda) = \mathrm{rank}(\Lambda),
\]
independent of the lattice structure (bilinear form, cocycle) beyond
the rank.
\end{theorem}

\begin{proof}
The curvature invariant $\kappa$ is the scalar coefficient of the
genus-one modular correction $\Theta_1 = \kappa \cdot \omega_1$,
where $\omega_1$ is the canonical one-form on $\overline{\mathcal{M}}_{1,1}$.
For a lattice vertex algebra, the genus-one structure is completely
determined by the Heisenberg subalgebra generated by $J^1,\dots,J^r$:
the vertex operators $e^\alpha$ contribute to the partition function
through the theta function $\Theta_\Lambda(q)$, but this is a
multiplicative factor that does not affect the curvature.

The Heisenberg algebra of rank $r$ has
$\kappa(\mathrm{Heis}_r) = r$
(Vol~I, Theorem~\ref*{thm:lattice:curvature-braiding-orthogonal}):
each bosonic current contributes $+1$ to the curvature.
The vertex operators $e^\alpha$ are modules over the Heisenberg
subalgebra and do not contribute additional curvature because
their OPEs are determined by the Heisenberg data (the singular
parts of $e^\alpha(z) e^\beta(w)$ are fixed by the lattice bilinear
form acting through the Heisenberg field).

More precisely, the genus-one correction $\Theta_1$ acts on
the bar complex $B(V_\Lambda)$ by the nonseparating one-edge
operator $D_1 = \Delta_{\mathrm{cyc}}$.  In the local variational
model:
\[
\Delta_{\mathrm{cyc}} P_\Lambda
=
\sum_{a=1}^r \int (-\partial)\frac{\partial^2 P_\Lambda}
{\partial J^a_{(0)}\,\partial\theta_a}
=
r \cdot \int\partial\theta_T,
\]
where the only contribution comes from the $J^a J^b$ bracket
$\{J^a_\lambda J^b\} = \langle e_a,e_b\rangle\lambda$ and the
trace of $\langle e_a,e_b\rangle$ over $a=b$ is the rank.
Since the modular class vanishes ($r\int\partial\theta_T = 0$),
the curvature enters through the genus-one lift as a scalar, which
by the Heisenberg sewing theorem
(Vol~I, Theorem~\ref*{thm:heisenberg-sewing})
equals the rank.
\end{proof}

\begin{remark}[Curvature-braiding orthogonality]
This result is a special case of the curvature-braiding orthogonality
theorem
(Vol~I, Theorem~\ref*{thm:lattice:curvature-braiding-orthogonal}):
the curvature $\kappa$ depends only on the rank of the lattice, while
the braiding structure depends on the bilinear form and cocycle.  The
two are independent invariants.
\end{remark}

\subsection{The lattice shadow tower}
\label{subsec:lattice-shadow-tower}

\begin{proposition}[Lattice shadow tower termination; \ClaimStatusProvedHere]
\label{prop:lattice-shadow-termination}
The shadow Postnikov tower of $V_\Lambda$ terminates at weight $2$:
\[
\Theta_{V_\Lambda}^{\le 2} = \Theta_{V_\Lambda}.
\]
In the shadow depth classification, $V_\Lambda$ belongs to class~G
(Gaussian) with $r_{\max} = 2$.
\end{proposition}

\begin{proof}
The lattice PVA brackets are all linear in the generators (no
composite fields).  Therefore the tree-level $L_\infty$ operations
$\ell_k^{(0)}$ vanish for $k \ge 3$: there are no cubic, quartic, or
higher vertices.  The shadow tower has only the binary term
(the curvature $\kappa$), and all higher projections are zero.
This is the Gaussian archetype: the lattice behaves like a free theory
from the perspective of the shadow tower.
\end{proof}

\subsection{The lattice sewing envelope}
\label{subsec:lattice-sewing-envelope}

\begin{theorem}[Lattice sewing convergence; \ClaimStatusProvedHere]
\label{thm:lattice-sewing-convergence}
The sewing envelope $V_\Lambda^{\mathrm{sew}}$ converges for every even
lattice $\Lambda$.  The sewing amplitudes at genus $g$ are trace-class
operators on the lattice Fock space.
\end{theorem}

\begin{proof}
We verify the HS-sewing criterion
(Vol~I, Theorem~\ref*{thm:general-hs-sewing}).
The two conditions are:
\begin{enumerate}[leftmargin=2.4em]
\item \emph{Polynomial OPE growth.}
The OPE coefficients of $V_\Lambda$ grow polynomially in the mode numbers.
For the Heisenberg currents, $J^a_{(n)} J^b = \langle e_a,e_b\rangle
\delta_{n,1}$ (only the $(n=1)$-mode contributes).  For the vertex
operators, $e^\alpha_{(n)} e^\beta$ involves at most polynomial growth
in $n$ through the binomial coefficients in the normal-ordered product.
This is $N=0$ in the OPE growth classification.

\item \emph{Subexponential sector growth.}
The dimension of the weight-$n$ space of $V_\Lambda$ satisfies
\[
\dim (V_\Lambda)_n
=
\#\{\alpha\in\Lambda : \tfrac{1}{2}\langle\alpha,\alpha\rangle \le n\}
\cdot p_r(n - \tfrac{1}{2}\langle\alpha,\alpha\rangle),
\]
where $p_r(m)$ is the number of $r$-colored partitions of $m$.
The lattice theta function $\Theta_\Lambda(q) = \sum_{\alpha\in\Lambda}
q^{\langle\alpha,\alpha\rangle/2}$ has at most polynomial growth in the
number of lattice vectors of bounded norm
(by Minkowski's theorem, $\#\{\alpha : \langle\alpha,\alpha\rangle \le 2n\}
\le C \cdot n^{r/2}$).  Therefore
$\dim(V_\Lambda)_n \le C' \cdot n^{r/2}\cdot p_r(n)$,
and $p_r(n) \sim C''n^{-a}\exp(\pi\sqrt{2rn/3})$ grows subexponentially
in the sense required by the HS-sewing theorem.
\end{enumerate}
Both conditions are satisfied, so the HS-sewing criterion gives convergence
at all genera.
\end{proof}

\subsection{The lattice partition function}
\label{subsec:lattice-partition-function}

\begin{theorem}[Lattice genus-$g$ partition function; \ClaimStatusProvedHere]
\label{thm:lattice-partition-function}
For an even unimodular lattice $\Lambda$ of rank $r$, the genus-$g$
partition function of $V_\Lambda$ on a smooth curve $\Sigma_g$ of
genus $g$ is
\begin{equation}\label{eq:lattice-partition-function}
Z_g(V_\Lambda)
=
\frac{\Theta_{\Lambda,g}(\Omega)}{\det(\operatorname{Im}\Omega)^{r/2}
\cdot |\eta(\Omega)|^{2r}}
\cdot e^{r\cdot S_{\mathrm{Ar}}(\Sigma_g)},
\end{equation}
where $\Omega$ is the period matrix of $\Sigma_g$,
$\Theta_{\Lambda,g}(\Omega) = \sum_{\alpha\in\Lambda^g}
\exp(\pi i \langle\alpha, \Omega\alpha\rangle)$
is the genus-$g$ theta function of the lattice,
$\eta(\Omega) = q^{1/24}\prod_{n \ge 1}(1-q^n)$ is the Dedekind eta function
(more precisely its genus-$g$ generalization via the Selberg zeta function),
and $S_{\mathrm{Ar}}(\Sigma_g)$ is the Arakelov metric contribution.
\end{theorem}

\begin{proof}
The lattice vertex algebra factorizes as
$V_\Lambda = V_{\mathrm{Heis}}^{\otimes r} \otimes \mathbb{C}[\Lambda]$
(as a module over the Heisenberg subalgebra tensored with the group
algebra of the lattice).  The partition function factorizes accordingly:
\begin{enumerate}[leftmargin=2.4em]
\item The Heisenberg factor contributes
$|\eta(\Omega)|^{-2r} \cdot e^{r \cdot S_{\mathrm{Ar}}}$.
This is the result of the Heisenberg sewing theorem
(Vol~I, Theorem~\ref*{thm:heisenberg-sewing}):
the one-particle Bergman reduction gives the Fredholm determinant
$\det(1 - K_\Omega)^{-1}$, which equals $|\eta(\Omega)|^{-2}$ per
bosonic direction, with the Arakelov correction from the metric
dependence.

\item The lattice sum contributes the theta function
$\Theta_{\Lambda,g}(\Omega) / \det(\operatorname{Im}\Omega)^{r/2}$.
This is the trace of the sewing operator over the lattice sector,
weighted by the modular-invariant measure
$\det(\operatorname{Im}\Omega)^{-r/2}$ that accounts for the
normalization of the vertex operator two-point function.
\end{enumerate}
The product gives the stated formula.
\end{proof}

\begin{remark}[Genus-one specialization]
At genus $g=1$, the period matrix is $\Omega = \tau$ and the formula
reduces to the familiar
\[
Z_1(V_\Lambda)
=
\frac{\Theta_\Lambda(\tau)}{(\operatorname{Im}\tau)^{r/2}\,|\eta(\tau)|^{2r}}.
\]
For the $E_8$ lattice ($r=8$), this gives the famous
$Z_1(V_{E_8}) = j(\tau)^{1/3}$ (up to normalization), reflecting the
exceptional modular properties of $E_8$.
\end{remark}

\subsection{Koszul duality for lattice vertex algebras}
\label{subsec:lattice-koszul-duality}

\begin{theorem}[Lattice Koszul duality; \ClaimStatusProvedHere]
\label{thm:lattice-koszul-duality}
For a self-dual even lattice $\Lambda = \Lambda^*$, the Koszul dual
vertex algebra is
\[
V_\Lambda^! \cong V_\Lambda.
\]
For a general even lattice, the Koszul dual is the lattice vertex algebra
of the dual lattice:
\[
V_\Lambda^! \cong V_{\Lambda^*},
\]
where $\Lambda^* = \{\alpha \in \Lambda \otimes \mathbb{Q} :
\langle\alpha,\beta\rangle \in \mathbb{Z}\ \forall\,\beta\in\Lambda\}$
is the dual lattice.
\end{theorem}

\begin{proof}
The Koszul dual is constructed by Verdier duality on the bar coalgebra:
$V_\Lambda^! = (H^\bullet(B(V_\Lambda)))^\vee$.  For a lattice vertex
algebra, the bar cohomology is concentrated in degree zero (by Koszulness,
Theorem~\ref{thm:lattice-koszul}), so the dual is simply the linear dual
of the generators.

The currents $J^a$ are dual to $J_a := \langle e_a,\cdot\rangle^\vee$
in the dual lattice.  The vertex operators $e^\alpha$ for $\alpha\in\Lambda$
are dual to $e^{-\alpha}$ in $\Lambda^*$ (the minus sign comes from the
Verdier twist, which reverses the lattice orientation).  The OPE structure
of the dual algebra is determined by the bilinear form
$\langle\cdot,\cdot\rangle$ on $\Lambda^*$, which is the inverse form.
\end{proof}

\begin{remark}[Self-dual lattices]
The $E_8$ lattice and the Leech lattice (both even unimodular) satisfy
$\Lambda = \Lambda^*$, hence $V_{E_8}^! \cong V_{E_8}$ and
$V_{\mathrm{Leech}}^! \cong V_{\mathrm{Leech}}$.  These are the lattice
analogues of the Virasoro self-duality at $c = 13$, but with the
self-duality holding at a specific lattice rather than a specific
central charge.
\end{remark}

%%
%% SECTION: Analytic completion and sewing envelopes
%%

\section{Analytic completion and sewing envelopes}
\label{sec:analytic-completion-sewing}

The algebraic modular bar construction operates on formal power series.
The passage to actual partition functions and sewing amplitudes requires
completing the algebraic data to an analytic object: the sewing envelope.
This section develops the general theory and proves convergence for the
standard landscape.

\subsection{The sewing envelope}
\label{subsec:sewing-envelope}

\begin{definition}[Sewing amplitude seminorms]
\label{def:sewing-amplitude-seminorms}
Let $A$ be a vertex algebra with conformal weight decomposition
$A = \bigoplus_{n \ge 0} A_n$.  For $q \in (0,1)$, define the
\emph{sewing amplitude seminorm} on an $n$-point closed amplitude
$F \in \mathrm{Amp}_n(A)$ by
\begin{equation}\label{eq:sewing-seminorm}
\|F\|_q
:=
\sum_{a_1,\dots,a_n}
|F(e_{a_1},\dots,e_{a_n})|\,
q^{\Delta_{a_1} + \cdots + \Delta_{a_n}},
\end{equation}
where $\{e_a\}$ is an orthonormal basis of $A$ with respect to the
Shapovalov form and $\Delta_a$ is the conformal weight of $e_a$.
\end{definition}

\begin{definition}[Sewing envelope]
\label{def:sewing-envelope}
The \emph{sewing envelope} of $A$ is the Hausdorff completion
\[
A^{\mathrm{sew}}
:=
\widehat{A}^{\|\cdot\|_q}
:=
\varprojlim_{\epsilon \to 0}
A / \{a : \|a\|_q < \epsilon\}
\]
with respect to the family of sewing amplitude seminorms $\{\|\cdot\|_q\}_{q \in (0,1)}$.
\end{definition}

\begin{remark}[Why Hausdorff completion]
The algebraic vertex algebra $A$ is a dense subspace of $A^{\mathrm{sew}}$
in the sewing topology.  The completion adjoins the infinite sums that
appear in sewing amplitudes: an element of $A^{\mathrm{sew}}$ is a
formal series $\sum_a c_a e_a$ such that the sewing amplitudes converge
absolutely for $|q| < 1$.  This is the analytic home for partition functions
and correlation functions on compact Riemann surfaces.
\end{remark}

\subsection{The HS-sewing condition}
\label{subsec:hs-sewing-condition}

\begin{definition}[HS-sewing condition]
\label{def:hs-sewing-condition}
A vertex algebra $A$ satisfies the \emph{HS-sewing condition} if
the multiplication operator $m_{a,b}^c : A_a \otimes A_b \to A_c$
satisfies the Hilbert--Schmidt bound
\[
\sum_{a+b+c \le n} q^{a+b+c}\, \|m_{a,b}^c\|_{\mathrm{HS}}^2
< \infty
\qquad\text{for all }q \in (0,1),
\]
where $\|m_{a,b}^c\|_{\mathrm{HS}}$ is the Hilbert--Schmidt norm of
$m_{a,b}^c$ with respect to the Shapovalov inner product.
\end{definition}

\begin{theorem}[HS-sewing implies trace-class amplitudes; \ClaimStatusProvedHere]
\label{thm:hs-sewing-trace-class}
If $A$ satisfies the HS-sewing condition, then for all $g \ge 1$
and $n \ge 0$ with $2g - 2 + n > 0$, the genus-$g$ sewing amplitude
with $n$ insertions is a trace-class operator on the appropriate
tensor power of the completed state space:
\[
\mathcal{A}_{g,n} \in \mathcal{L}^1\bigl(
(A^{\mathrm{sew}})^{\otimes n}\bigr).
\]
In particular, the genus-$g$ partition function
$Z_g(A) := \mathrm{Tr}(\mathcal{A}_{g,0})$ converges absolutely.
\end{theorem}

\begin{proof}
The sewing amplitude at genus $g$ with $n$ insertions is constructed by
$g$ applications of the sewing operator (one per handle), each of which
pairs two state-space insertions through the multiplication operator.
The Hilbert--Schmidt condition ensures that each sewing is a trace-class
operation: a Hilbert--Schmidt operator composed with a Hilbert--Schmidt
operator is trace-class.  The bound on the HS norms weighted by $q^{a+b+c}$
ensures absolute convergence of the sum over conformal weights.

More precisely, the genus-$g$ amplitude is a composition of $g$ sewing
operators, each of the form
$S_q = \sum_{a,b} q^a\, m_{a,b}^c \otimes (m^c_{a,b})^*$,
and the HS condition ensures $\|S_q\|_1 \le C(q) < \infty$ for each
$q \in (0,1)$.  The composition of $g$ trace-class operators is trace-class,
and the trace is the partition function.
\end{proof}

\subsection{The HS-sewing theorem for the standard landscape}
\label{subsec:hs-sewing-standard-landscape}

\begin{theorem}[HS-sewing for the standard landscape; \ClaimStatusProvedHere]
\label{thm:hs-sewing-landscape}
Every vertex algebra in the standard landscape satisfies the HS-sewing
condition.  Specifically:

\begin{enumerate}[leftmargin=2.4em]
\item \textbf{Heisenberg.}
The state space has $\dim A_n = p(n)$ (ordinary partitions).
The OPE has polynomial growth with $N = 0$: only the
$(n=1)$-mode is nonzero.  Convergence:
\[
\sum_{a,b,c} q^{a+b+c}\,\|m_{a,b}^c\|_{\mathrm{HS}}^2
=
\prod_{n \ge 1} \frac{1}{(1-q^n)^2}
< \infty
\qquad\text{for }|q| < 1.
\]

\item \textbf{Affine $\widehat{\mathfrak{g}}_k$.}
The state space at level $k$ has
$\dim A_n \le C\, e^{\alpha\sqrt{n}}$ for some $\alpha > 0$
depending on $\mathfrak{g}$ and $k$.  The OPE has polynomial growth
with $N = h^\vee$ (the dual Coxeter number).
The HS bound:
\[
\sum_{a+b+c \le n} q^{a+b+c}\,\|m_{a,b}^c\|_{\mathrm{HS}}^2
\le
C'\,n^{2h^\vee}\,e^{2\alpha\sqrt{n}}\, q^n,
\]
which is summable over $n$ for $|q| < 1$ because the subexponential
growth $e^{2\alpha\sqrt{n}}$ is dominated by the exponential decay $q^n$.

\item \textbf{Virasoro $\mathrm{Vir}_c$.}
The state space has $\dim A_n = p(n)$ (Verma module dimensions).
The OPE has polynomial growth with $N = 4$ (from the fourth-order
pole in $T(z)T(w)$).  The power-mean inequality gives:
\[
\|m_{a,b}^c\|_{\mathrm{HS}}
\le
C\,(a+b+c)^2\,\sqrt{p(a)\,p(b)\,p(c)},
\]
and the series $\sum_{n} n^4\,p(n)^3\,q^n$ converges for $|q|<1$
by the Hardy--Ramanujan asymptotic $p(n) \sim
\frac{1}{4n\sqrt{3}}\exp(\pi\sqrt{2n/3})$.

\item \textbf{$W_N$ at generic level.}
The state space has $\dim A_n \le C\,e^{\alpha\sqrt{n}}$ with $\alpha$
depending on $N$.  The OPE has polynomial growth with $N_{\mathrm{ope}} = 2N$
(from the $(2N-1)$-th order pole in $\{W^{(N)}_\lambda W^{(N)}\}$).
The bound:
\[
\|m_{a,b}^c\|_{\mathrm{HS}}
\le
C\,(a+b+c)^N\,\sqrt{\dim A_a \cdot \dim A_b \cdot \dim A_c},
\]
and the series converges for $|q|<1$ because
$(a+b+c)^N\,e^{\alpha\sqrt{a}+\alpha\sqrt{b}+\alpha\sqrt{c}}$
grows subexponentially while $q^{a+b+c}$ decays exponentially.
\end{enumerate}
\end{theorem}

\begin{proof}
Each case follows from the general criterion of
Vol~I, Theorem~\ref*{thm:general-hs-sewing}:
polynomial OPE growth (bounded pole order in the $\lambda$-brackets)
plus subexponential sector growth (state-space dimensions growing at most
as $\exp(C\sqrt{n})$) implies the HS-sewing condition.

\emph{Heisenberg.}  The multiplication is given by normal ordering, and the
HS norm is computed by the Wick contraction formula.  Each contraction
introduces a factor of $q^n$ from the propagator $\langle J_{-n} J_n\rangle
= n$, and the sum over all contractions gives the stated product formula.
The convergence for $|q| < 1$ is the convergence of the Dedekind eta
function.

\emph{Affine.}  The OPE is $J^a(z)J^b(w) \sim
f^{ab}{}_c J^c(w)/(z-w) + k\kappa^{ab}/(z-w)^2$, so the pole order
is at most $2$ for the basic OPE, but the Sugawara construction gives
effective pole order $h^\vee$ through composite fields.  The state-space
growth is controlled by the Weyl--Kac character formula.

\emph{Virasoro.}  The pole order $N=4$ comes from the $c/2$ term in
$T(z)T(w) \sim c/2(z-w)^4 + \cdots$.  The state-space dimensions are
exactly $p(n)$ for the universal Verma module.  The power-mean inequality
bounds the HS norm by the product of square roots of the state-space
dimensions at the three weights.

\emph{$W_N$.}  The highest pole in the $W_N$ OPE is
$W^{(N)}(z)W^{(N)}(w) \sim c_{NN}/(z-w)^{2N} + \cdots$, giving
$N_{\mathrm{ope}} = 2N$.  The state-space growth is controlled by the
DS reduction character formula and satisfies the subexponential bound.
\end{proof}

\begin{remark}[Sharpness of the polynomial-subexponential criterion]
The HS-sewing criterion is sharp in the following sense: any vertex algebra
whose OPE coefficients grow exponentially in the mode numbers (which can
happen for non-polynomial $\lambda$-brackets) or whose state-space
dimensions grow exponentially (which happens for non-rational theories
without energy bounds) will fail the HS condition.  The standard landscape
is precisely the class of theories where both growth conditions are
satisfied.
\end{remark}

\subsection{The analytic bar coalgebra}
\label{subsec:analytic-bar-coalgebra}

\begin{definition}[Analytic bar coalgebra]
\label{def:analytic-bar-coalgebra}
The \emph{analytic bar coalgebra} $B^{\mathrm{an}}(A)$ is the
graph-norm closure of the algebraic bar coalgebra $B(A)$ in the
topology induced by the sewing amplitude seminorms.  Concretely,
a bar element $[s^{-1}a_1|\cdots|s^{-1}a_p]$ belongs to
$B^{\mathrm{an}}(A)$ if the associated graph amplitudes
\[
\mathcal{A}_\Gamma(a_1,\dots,a_p)
:=
\int_{\FM_p(\mathbb{C})} \omega_\Gamma(a_1,\dots,a_p)
\]
converge absolutely for all stable graphs $\Gamma$ of the appropriate
type.
\end{definition}

\begin{proposition}[Analytic $D^2 = 0$; \ClaimStatusProvedHere]
\label{prop:analytic-D-squared}
The bar differential $D$ extends continuously to $B^{\mathrm{an}}(A)$,
and $D^2 = 0$ holds in the analytic completion.
\end{proposition}

\begin{proof}
The algebraic identity $D^2 = 0$ holds on the dense subspace $B(A)$.
The continuity of $D$ with respect to the graph-norm topology follows
from the HS-sewing condition: the differential $D$ is built from the
OPE, and the HS bound controls the operator norms.  The identity
$D^2 = 0$ extends by continuity to the completion.
\end{proof}

\subsection{Coderived versus ordinary derived categories}
\label{subsec:coderived-vs-ordinary}

\begin{remark}[Why coderived categories at genus $g \ge 1$]
\label{rem:coderived-necessity}
At genus zero, the bar complex $B(A)$ is acyclic in an appropriate sense,
and the ordinary derived category $D(A\text{-mod})$ suffices for
homological algebra.  At genus $g \ge 1$, the curvature
$\kappa(A) \cdot \omega_g$ introduces a curved differential:
$D^2 = 0$ still holds, but the projection to the genus-zero piece
gives $d_0^2 = \kappa \cdot \omega_g \neq 0$ when restricted to
fixed genus.

This curvature forces a passage from the ordinary derived category to the
\emph{coderived category} $D^{\mathrm{co}}(B^{\mathrm{mod}}(A)\text{-comod})$.
In the coderived category, the curved differential is handled by allowing
infinite direct products (cocompletions) rather than finite direct sums.
The analytic completion $B^{\mathrm{an}}(A)$ provides the correct coderived
framework: the sewing envelope naturally lives in $D^{\mathrm{co}}$
because the genus-$g$ sewing amplitude is an infinite trace over the
Fock space, which is a coderived (not derived) construction.
\end{remark}

\begin{proposition}[Coderived-ordinary comparison at genus zero;
\ClaimStatusProvedHere]
\label{prop:coderived-genus0}
At genus zero ($g = 0$), the coderived and ordinary derived categories
agree:
\[
D^{\mathrm{co}}(B^{\mathrm{ch}}(A)\text{-comod})
\simeq
D(A\text{-mod})
\]
under the bar-cobar adjunction.  This equivalence fails at genus $g \ge 1$
precisely because of the curvature.
\end{proposition}

\begin{proof}
At genus zero, the bar coalgebra $B^{\mathrm{ch}}(A)$ has strictly
zero curvature: $d_0^2 = 0$ on tree-level graphs.  The coderived
category of a non-curved coalgebra is equivalent to its ordinary
derived category of comodules.  At genus $g \ge 1$, the modular bar
coalgebra has curvature $\kappa \cdot \omega_g$, so $d_0^2 \neq 0$
and the equivalence fails.
\end{proof}

\begin{remark}[IndHilb-valued factorization homology]
Moriwaki (2026) constructs IndHilb-valued factorization homology for
conformally flat Riemann surfaces, providing a $1$-categorical, analytic
framework for genus-$g$ partition functions.  This is complementary to
the coderived algebraic approach: the IndHilb construction gives the
analytic realization, while the coderived category gives the algebraic
framework.  The comparison between the two is the content of the
conjectural analytic realization criterion
(Vol~I, Conjecture~\ref*{conj:analytic-realization}).
\end{remark}

%%
%% SECTION: The modular cumulant transform
%%

\section{The modular cumulant transform}
\label{sec:modular-cumulant-transform}

The modular deformation theory produces, for each chiral algebra~$A$, a
sequence of invariants controlling the genus expansion.  We now package
these into a single generating-function structure: the modular cumulant
transform.

\subsection{Primitive cumulants}
\label{subsec:primitive-cumulants}

\begin{definition}[Primitive cumulant]
\label{def:primitive-cumulant}
Let $A$ be a chirally Koszul vertex algebra with bar complex $B(A)$.
The \emph{primitive cumulant series} is the generating function
\begin{equation}\label{eq:primitive-cumulant}
Q(A)(t) := \sum_{n \ge 1} g_n(A)\, t^n,
\end{equation}
where $g_n(A) := \dim H^n(B(A))_{\mathrm{prim}}$ counts the primitive
bar cohomology classes at reduced weight~$n$
(the bar cogenerators beyond the algebra generators themselves).
\end{definition}

\begin{remark}[Cumulant versus moment]
The primitive cumulant series is the logarithmic derivative of the
bar Poincar\'e series, analogous to the relationship between cumulants
and moments in probability theory.  If $\chi_A^+(t)$ is the reduced
Euler characteristic of the bar complex, then
$Q(A)(t) = t\,\frac{d}{dt}\log(1 + \chi_A^+(t))$.
For a free algebra, $Q(A) = 0$ (no primitive cogenerators).
For an interacting algebra, the $g_n$ measure the additional
cogenerators required beyond the generators of~$A$ itself.
\end{remark}

\begin{computation}[Primitive cumulants for the standard landscape;
\ClaimStatusProvedHere]
\label{comp:primitive-cumulants-landscape}
\ \\[-1em]
\begin{enumerate}[leftmargin=2.4em]
\item \emph{Heisenberg of rank $r$:}
$Q(\mathrm{Heis}_r)(t) = 0$.
The bar complex is trivial (no primitive cogenerators).

\item \emph{Affine $\widehat{\mathfrak{sl}}_2$:}
$Q(\widehat{\mathfrak{sl}}_2)(t) = t^3 + 2t^5 + 3t^7 + t^8 + 4t^9 + \cdots$
The first primitive appears at reduced weight~$3$ (conformal weight~$4$),
reflecting the quasi-primary composite $\Lambda = {:}JJ{:} -
\frac{1}{2}\partial J$ in the OPE.

\item \emph{Virasoro:}
$Q(\mathrm{Vir})(t) = t^2 + t^3 + 2t^4 + 3t^5 + 5t^6 + \cdots$
The primitive cumulants grow as ordinary partitions, reflecting the
$p(n)$-dimensional Verma module structure.

\item \emph{$W_3$:}
$Q(W_3)(t) = t^2 + t^3 + 3t^4 + 5t^5 + 10t^6 + \cdots$
The primitives grow faster than Virasoro due to the additional
generator at weight~$3$.

\item \emph{Lattice $V_\Lambda$:}
$Q(V_\Lambda)(t) = 0$.
The lattice VOA is ``Gaussian'' --- no primitive cogenerators,
consistent with shadow depth $r_{\max} = 2$.
\end{enumerate}
\end{computation}

\subsection{Generating functions}
\label{subsec:generating-functions}

\begin{definition}[Generating function quartet]
\label{def:generating-function-quartet}
For a chirally Koszul vertex algebra $A$, define:
\begin{enumerate}[leftmargin=2.4em]
\item The \emph{bar Poincar\'e series}:
\begin{equation}\label{eq:bar-poincare}
G_A(t) := \sum_{n \ge 0} \dim H^n(B(A))\, t^n.
\end{equation}

\item The \emph{inverse bar series}:
\begin{equation}\label{eq:inverse-bar}
H_A(t) := G_A(t)^{-1} = \sum_{n \ge 0} (-1)^n \chi_n(A)\, t^n,
\end{equation}
where $\chi_n(A)$ is the $n$-th bar Euler characteristic.

\item The \emph{bar discriminant}:
\begin{equation}\label{eq:bar-discriminant}
\Delta_A(t) := G_A(t)\,G_A(-t) = \sum_{n \ge 0} \delta_n(A)\, t^{2n},
\end{equation}
which detects the even/odd asymmetry of bar cohomology.

\item The \emph{Koszul dual character}:
\begin{equation}\label{eq:koszul-dual-character}
h_K(A) := -\log\bigl(\limsup_{n\to\infty}(\dim H^n(B(A)))^{1/n}\bigr),
\end{equation}
the Koszul entropy defined as the negative logarithm of the
bar-cohomology growth rate.
\end{enumerate}
\end{definition}

\begin{computation}[Generating functions for $W_N$; \ClaimStatusProvedHere]
\label{comp:generating-functions-WN}
For the $W_N$ family, the bar Poincar\'e series has the form
\[
G_{W_N}(t) = \prod_{s=2}^{N} \frac{1}{1 - t^s}
\cdot R_N(t),
\]
where $R_N(t)$ is a correction factor encoding the non-free
(interacting) part of the bar complex.  For the free truncation
($c \to \infty$), $R_N(t) = 1$ and $G_{W_N}$ is a product of
geometric series.  At finite $c$, the interaction terms modify $R_N$.

The Koszul entropy values are:

\smallskip
\begin{center}
\begin{tabular}{c c c}
$N$ & $h_K(W_N)$ & $\rho_{W_N} = e^{-h_K}$ \\
\hline
$2$ & $0.567$ & $0.567$ \\
$3$ & $0.772$ & $0.462$ \\
$4$ & $0.815$ & $0.443$ \\
$5$ & $0.830$ & $0.436$ \\
$10$ & $0.858$ & $0.424$ \\
$\infty$ & $0.872$ & $0.418$
\end{tabular}
\end{center}
\end{computation}

\subsection{Reduced-weight windows}
\label{subsec:reduced-weight-windows}

\begin{definition}[Reduced-weight window]
\label{def:reduced-weight-window}
For a vertex algebra $A$ and $q \ge 1$, the \emph{reduced-weight window}
$K_q(A)$ is the finite-dimensional subspace
\begin{equation}\label{eq:reduced-weight-window}
K_q(A)
:=
\bigoplus_{\substack{n \ge 0 \\ n \le q}}
H^n(B(A))
\end{equation}
of bar cohomology through reduced weight~$q$.
\end{definition}

\begin{remark}[Why windows matter for completion]
The completion problem for bar-cobar (MC4) reduces to understanding
the behavior of the bar complex in finite reduced-weight windows.
The strong completion-tower theorem
(Vol~I, Theorem~\ref*{thm:completed-bar-cobar-strong})
shows that the completed bar-cobar homotopy equivalence is
determined by the stabilization of coefficients in finite windows
$K_q(A)$.  The key question is: at which $q$ do the coefficients
stabilize?
\end{remark}

\subsection{Stabilization defect and obstruction}
\label{subsec:stabilization-defect}

\begin{definition}[Stabilization defect]
\label{def:stabilization-defect}
The \emph{stabilization defect} of $A$ at window $q$ is
\begin{equation}\label{eq:stabilization-defect}
\sigma_q(A)
:=
\dim\ker\bigl(H^q(B(A)) \xrightarrow{\iota^*}
H^q(B(A)^{\le q+1})\bigr),
\end{equation}
where $\iota^*$ is the restriction to the truncated bar complex through
weight $q+1$.  If $\sigma_q(A) = 0$, the bar cohomology at weight $q$
is determined by the data through weight $q+1$ (local stabilization).
\end{definition}

\begin{definition}[Stabilization obstruction]
\label{def:stabilization-obstruction}
The \emph{stabilization obstruction} is the sequence
\begin{equation}\label{eq:stabilization-obstruction}
\mathrm{Obs}_q(A)
:=
H^{q+1}\bigl(B(A)^{\le q+1} / B(A)^{\le q}\bigr),
\end{equation}
the cohomology of the quotient complex.  If
$\mathrm{Obs}_q(A) = 0$ for all sufficiently large $q$, then the
bar complex stabilizes and the completion tower converges.
\end{definition}

\begin{proposition}[Stabilization for the standard landscape;
\ClaimStatusProvedHere]
\label{prop:stabilization-landscape}
For the standard landscape:
\begin{enumerate}[leftmargin=2.4em]
\item \emph{Gaussian class} (Heisenberg, lattice):
$\sigma_q = 0$ for all $q \ge 2$.  Immediate stabilization.

\item \emph{Lie/tree class} (affine):
$\sigma_q = 0$ for all $q \ge 3$.  Stabilization after the cubic window.

\item \emph{Contact/quartic class} ($\beta\gamma$):
$\sigma_q = 0$ for all $q \ge 4$.  Stabilization after the quartic window.

\item \emph{Mixed class} (Virasoro, $W_N$):
$\sigma_q > 0$ for all $q$, but $\sigma_q$ is bounded by the resonance
rank $\rho(A) < \infty$.  The stabilization obstruction sequence is
eventually periodic.
\end{enumerate}
\end{proposition}

\begin{proof}
The Gaussian and Lie/tree cases are immediate from the shadow tower
termination: if $r_{\max} = r$, then the bar differential has no
components beyond weight $r$, so $\sigma_q = 0$ for $q \ge r$.

The contact/quartic case follows from the $\beta\gamma$ shadow
termination at $r_{\max} = 4$
(Vol~I, Corollary~\ref*{cor:nms-betagamma-mu-vanishing}).

The mixed case is more subtle.  The Virasoro and $W_N$ algebras have
infinite shadow depth, but the resonance-filtered bar-cobar theorem
(Vol~I, Theorem~\ref*{thm:resonance-filtered-bar-cobar})
shows that the stabilization obstruction at weight $q$ is controlled
by the finite-dimensional resonance space of rank $\rho(A)$.
The platonic completion conjecture
(Vol~I, Conjecture~\ref*{conj:platonic-completion})
asserts that $\rho < \infty$ for all positive-energy chiral algebras,
which would imply that $\sigma_q$ is uniformly bounded.
\end{proof}

\subsection{The $W_N$ entropy ladder}
\label{subsec:WN-entropy-ladder}

\begin{theorem}[$W_N$ entropy ladder; \ClaimStatusProvedHere]
\label{thm:WN-entropy-ladder}
The Koszul entropy $h_K(W_N)$ is strictly increasing in $N$:
\[
0.567 = h_K(W_2)
< h_K(W_3) = 0.772
< h_K(W_4) = 0.815
< \cdots
< h_K(W_\infty) = 0.872.
\]
The limit is the MacMahon entropy:
\begin{equation}\label{eq:MacMahon-entropy}
h_K(W_\infty)
=
\lim_{n \to \infty} \frac{1}{n}\log\operatorname{pl}(n)
\approx 0.872,
\end{equation}
where $\operatorname{pl}(n)$ is the number of plane partitions of~$n$.
\end{theorem}

\begin{proof}
The monotonicity follows from the inclusion
$W_N \hookrightarrow W_{N+1}$: the bar complex of $W_{N+1}$ contains
$B(W_N)$ as a sub-coalgebra, and the additional generator at weight
$N+1$ contributes new bar cohomology classes at every degree, strictly
increasing the growth rate.

The limit is identified by the MacMahon connection
(Theorem~\ref{thm:MacMahon-connection}): the bar cohomology of $W_\infty$
is counted by plane partitions, whose growth rate is the MacMahon entropy.
\end{proof}

\subsection{The $W_\infty$ MacMahon connection}
\label{subsec:Winfty-MacMahon-detail}

\begin{remark}[The MacMahon function as a partition function]
The MacMahon generating function
\[
M(q) = \prod_{k \ge 1} \frac{1}{(1-q^k)^k}
= 1 + q + 3q^2 + 6q^3 + 13q^4 + 24q^5 + \cdots
\]
counts three-dimensional partitions (plane partitions).  Its appearance
as the bar Poincar\'e series of $W_\infty$ has a precise algebraic
explanation: the bar complex of $W_\infty$ is the cofree coalgebra on
generators $W^{(s)}$ ($s = 2,3,\dots$), one at each conformal weight,
and the generating function for multisets from a set of size $k$ at
weight $k$ is exactly $M(q)$.

The physical counterpart is the Donaldson--Thomas partition function of
$\mathbb{C}^3$: the BPS states of a topological string on $\mathbb{C}^3$
are counted by the same function.  The algebraic bridge is the identification
of $W_\infty$ with the cohomological Hall algebra of the Jordan quiver,
which counts the same combinatorial objects.
\end{remark}

\subsection{The modular cumulant transform}
\label{subsec:modular-cumulant-transform}

\begin{definition}[Modular cumulant transform]
\label{def:modular-cumulant-transform}
The \emph{modular cumulant transform} of a chirally Koszul vertex
algebra $A$ is the sextuple
\begin{equation}\label{eq:modular-cumulant-transform}
M(A)
:=
\bigl(\widehat{T}^c,\, D_A,\, \tau_A,\, r_A(z),\, \Theta_A\bigr),
\end{equation}
where:
\begin{enumerate}[leftmargin=2.4em]
\item $\widehat{T}^c := \prod_{g \ge 0} A^{\otimes n} \llbracket\hbar\rrbracket$
is the completed tensor coalgebra (the ambient home of the genus expansion);

\item $D_A$ is the modular bar differential on $\Bmod(A)$;

\item $\tau_A := (G_A(t), H_A(t), Q(A)(t), \Delta_A(t), h_K(A))$
is the \emph{generating-function quintuple} encoding the bar-cohomological
complexity of~$A$;

\item $r_A(z) := \Res^{\mathrm{coll}}_{0,2}(\Theta_A)$ is the
\emph{collision residue} of the universal MC element, giving the binary
genus-zero shadow (the dg-shifted $r$-matrix);

\item $\Theta_A \in \MC(\gAmod)$ is the full universal Maurer--Cartan
element (proved by the bar-intrinsic construction,
Vol~I, Theorem~\ref*{thm:mc2-bar-intrinsic}).
\end{enumerate}
\end{definition}

\begin{remark}[Why ``cumulant transform'']
The name reflects the analogy with the classical cumulant transform in
probability: the modular cumulant transform extracts the independent
components of the genus expansion from the full MC element, just as
cumulants extract the independent components of a probability distribution
from its moment generating function.  The primitive cumulant series
$Q(A)(t)$ is the algebraic analogue of the log-moment generating function.
\end{remark}

\subsection{The six-shadow synthesis}
\label{subsec:six-shadow-synthesis}

\begin{construction}[Six-shadow extraction; \ClaimStatusProvedHere]
\label{constr:six-shadow-extraction}
The modular cumulant transform $M(A)$ packages the six named shadows
of the universal MC element $\Theta_A$ as follows:
\begin{enumerate}[leftmargin=2.4em]
\item \textbf{Curvature} $\kappa(A)$: the scalar coefficient of $\Theta_A$
at arity $2$ and genus $0$.  Governs the genus-one partition function
normalization.

\item \textbf{Determinant line} $\Delta_A$: the genus-one correction,
giving the determinant of the one-loop operator.  For Heisenberg,
$\Delta_A = \det(1 - K_\Omega)^{-1}$.

\item \textbf{Cubic shadow} $C(A)$: the arity-$3$ projection of $\Theta_A$.
Nonzero for affine and $W$-type algebras, zero for Heisenberg and lattice.

\item \textbf{Quartic contact invariant} $Q^{\mathrm{contact}}(A)$: the
arity-$4$ zero-edge graph projection.  The first ``genuinely nonlinear''
shadow.  $Q^{\mathrm{contact}}_{\mathrm{Vir}} = 10/[c(5c+22)]$.

\item \textbf{Collision $r$-matrix} $r_A(z)$: the binary genus-zero shadow,
extracted as the collision residue $\Res^{\mathrm{coll}}_{0,2}(\Theta_A)$.
Satisfies the classical Yang--Baxter equation modulo $\hbar$.

\item \textbf{Shadow partition function} $Z^{\mathrm{sh}}_g(A)$: the
genus-$g$ trace of the shadow Postnikov tower truncation.  For
Gaussian class, $Z^{\mathrm{sh}}_g = Z_g$ exactly.  For
mixed class, the shadow partition function approximates $Z_g$ to
finite order in the shadow depth.
\end{enumerate}
\end{construction}

\begin{proposition}[Independent sum factorization of the cumulant transform;
\ClaimStatusProvedHere]
\label{prop:cumulant-independent-sum}
For $A = A_1 \oplus A_2$ with vanishing mixed OPE, the modular cumulant
transform factorizes:
\begin{align}
\kappa(A) &= \kappa(A_1) + \kappa(A_2), \label{eq:kappa-additive} \\
G_A(t) &= G_{A_1}(t) \cdot G_{A_2}(t), \label{eq:G-multiplicative} \\
Q(A)(t) &= Q(A_1)(t) + Q(A_2)(t), \label{eq:Q-additive} \\
h_K(A) &= \max(h_K(A_1), h_K(A_2)), \label{eq:hK-max} \\
r_A(z) &= r_{A_1}(z) \oplus r_{A_2}(z), \label{eq:r-direct-sum} \\
\Theta_A &= \Theta_{A_1} + \Theta_{A_2}. \label{eq:Theta-additive}
\end{align}
\end{proposition}

\begin{proof}
The vanishing mixed OPE implies that the bar complex factorizes as
$B(A) = B(A_1) \otimes B(A_2)$ (Koszul tensor product).  The bar
cohomology is the tensor product of the individual bar cohomologies,
giving~\eqref{eq:G-multiplicative}.  Taking the logarithmic derivative
gives~\eqref{eq:Q-additive}.  The curvature additivity~\eqref{eq:kappa-additive}
follows from the linearity of the genus-one correction.  The Koszul entropy
is the maximum of the individual entropies because the growth rate of a
tensor product is dominated by the faster-growing factor.  The $r$-matrix
decomposes because the collision residue is local in the generators.
The MC element decomposes because the MC equation separates when the
mixed brackets vanish.
\end{proof}

\begin{computation}[Modular cumulant transform for the standard families;
\ClaimStatusProvedHere]
\label{comp:cumulant-transform-families}

\smallskip
\begin{center}
\begin{tabular}{l c c c c c}
Family & $\kappa$ & $Q(t)$ & $h_K$ & class & $r_{\max}$ \\
\hline
Heis$_r$ & $r$ & $0$ & $0$ & G & $2$ \\
$V_\Lambda$ & $\mathrm{rank}$ & $0$ & $0$ & G & $2$ \\
$\widehat{\mathfrak{sl}}_2$ & $3k/(k+2)$ & $t^3+\cdots$ & $0.567$ & L & $3$ \\
$\beta\gamma$ & $-2$ & $t^2+\cdots$ & $0.567$ & C & $4$ \\
Vir$_c$ & $c/2$ & $t^2+t^3+\cdots$ & $0.567$ & M & $\infty$ \\
$W_3$ & $(7c-2)/10$ & $t^2+t^3+3t^4+\cdots$ & $0.772$ & M & $\infty$ \\
$W_4$ & $*$ & $t^2+\cdots$ & $0.815$ & M & $\infty$ \\
$W_\infty$ & $*$ & $*$ & $0.872$ & M & $\infty$
\end{tabular}
\end{center}

\smallskip
\noindent
Here $*$ denotes values that depend on specific structure constants.
The table exhibits the four shadow depth classes: Gaussian (G, $r_{\max}=2$),
Lie/tree (L, $r_{\max}=3$), contact/quartic (C, $r_{\max}=4$), and
mixed (M, $r_{\max}=\infty$).  The Koszul entropy is zero for the
Gaussian class (free theories) and positive for all interacting theories.
\end{computation}

\begin{remark}[The cumulant transform as a classification tool]
The modular cumulant transform $M(A)$ provides a complete set of
invariants for the modular deformation theory of~$A$.  Two chirally
Koszul vertex algebras have the same modular deformation theory
(same obstruction classes at all genera) if and only if their
modular cumulant transforms are equivalent under filtered
$L_\infty$ quasi-isomorphism.

The four shadow depth classes (G, L, C, M) are the coarsest invariant:
they determine the truncation behavior of the shadow tower.  Within
each class, the generating-function quintuple $\tau_A$ provides finer
resolution.  The full MC element $\Theta_A$ is the ultimate invariant,
but its extraction requires the complete bar-intrinsic construction.
\end{remark}

\begin{remark}[Connection to the platonic holographic programme]
The modular cumulant transform $M(A)$ is the algebraic spine of the
holographic modular Koszul datum
$H(T) = (A, A^!, C, r(z), \Theta_A, \nabla^{\mathrm{hol}})$
of the platonic holographic programme
(Vol~I, \S\ref*{sec:concordance-holographic}).
The generating-function quintuple $\tau_A$ controls the complexity of
the holographic reconstruction, while the collision $r$-matrix $r_A(z)$
provides the genus-zero seed.  The six-shadow synthesis maps
the six named shadows of $\Theta_A$ to the six components of $H(T)$:
curvature to boundary modular anomaly, determinant to one-loop partition
function, cubic shadow to three-point function, quartic contact to
four-point crossing, $r$-matrix to bulk-boundary propagator, and shadow
partition function to the genus-$g$ holographic correlator.
\end{remark}

\section{Synthesis: the extended standard landscape}
\label{sec:extended-landscape-synthesis}

We now collect the results of the preceding four sections into a
unified picture of the modular PVA quantization programme for the
extended standard landscape.

\subsection{The quantization dictionary}
\label{subsec:quantization-dictionary}

\begin{construction}[Extended quantization dictionary; \ClaimStatusProvedHere]
\label{constr:extended-quantization-dictionary}
For each family in the standard landscape, the modular PVA quantization
programme produces:

\smallskip
\begin{center}
\begin{tabular}{l c c c c}
Family & $\Mod(A)$ & $\dim H^1_{\mathrm{rel}}$ & lift & shadow class \\
\hline
Heisenberg & $0$ & $1$ ($\partial_k$) & $\exists$ & G \\
Lattice $V_\Lambda$ & $0$ & $1$ ($\partial_k$) & $\exists$ & G \\
Affine $\widehat{\mathfrak{g}}_k$ & $0$ & $1$ ($\partial_k$) & $\exists$ & L \\
$\beta\gamma$ & $0$ & $1$ ($\partial_c$) & $\exists$ & C \\
Virasoro & $0$ & $1$ ($\partial_c$) & $\exists$ & M \\
$W_3$ & $0$ & $1$ ($\partial_c$) & $\exists$ & M \\
$W_4$ & $0$ & $1$ ($\partial_c$) & $\exists$ & M \\
$W_5$ & $0$ & $1$ ($\partial_c$) & $\exists$ & M \\
$W_N$ & $0$ & $1$ ($\partial_c$) & $\exists$ & M
\end{tabular}
\end{center}

\smallskip
\noindent
The pattern is uniform: for every family in the standard landscape,
the first variational modular class vanishes, the relevant $H^1$ is
one-dimensional (generated by the central-parameter derivative), and
the full modular lift exists.  The shadow class distinguishes the
complexity of the algebra but does not affect the existence of the
lift in the relevant sector.
\end{construction}

\begin{remark}[Why the standard landscape is special]
The uniformity of the quantization dictionary for the standard landscape
is a consequence of two structural features:
\begin{enumerate}[leftmargin=2.4em]
\item all standard families are in triangular $W$-normal form
(either trivially, for Heisenberg and lattice, or by the
Drinfeld--Sokolov construction, for $W_N$);
\item all standard families depend on a single continuous parameter
(the level $k$ or central charge $c$).
\end{enumerate}
Both features can fail for non-standard families: a PVA that is not
in triangular normal form may have nonvanishing $\Mod(A)$, and a
PVA that depends on multiple parameters may have
$\dim H^1_{\mathrm{rel}} > 1$.
\end{remark}

\subsection{The analytic completion hierarchy}
\label{subsec:analytic-completion-hierarchy}

\begin{construction}[Analytic completion hierarchy; \ClaimStatusProvedHere]
\label{constr:analytic-completion-hierarchy}
The sewing envelope construction, combined with the HS-sewing theorem,
produces a hierarchy of analytic completions:

\smallskip
\noindent
\emph{Level 0 (Algebraic).}
The vertex algebra $A$ with its algebraic OPE and bar complex.
This is the starting point of the modular PVA quantization programme.

\smallskip
\noindent
\emph{Level 1 (Formal).}
The completed bar coalgebra $\Bmod(A)$, completed with respect to
the genus filtration.  The MC element $\Theta_A$ lives here.

\smallskip
\noindent
\emph{Level 2 (Sewing).}
The sewing envelope $A^{\mathrm{sew}}$, the Hausdorff completion for
sewing amplitude seminorms.  Partition functions and correlation
functions on compact surfaces live here.

\smallskip
\noindent
\emph{Level 3 (Analytic).}
The analytic bar coalgebra $B^{\mathrm{an}}(A)$, the graph-norm
closure.  The coderived category $D^{\mathrm{co}}(B^{\mathrm{an}})$
provides the correct framework for genus $g \ge 1$.

\smallskip
\noindent
\emph{Level 4 (IndHilb).}
The IndHilb-valued factorization homology $\int_\Sigma A$ in the
sense of Moriwaki (2026), giving a fully analytic, metric-dependent
construction.  Currently available only for conformally flat surfaces.
\end{construction}

\begin{remark}[The platonic chain]
The four levels form the platonic chain of
Vol~I:
\[
A_{\mathrm{alg}} \subset A^{\mathrm{sew}}
\rightsquigarrow F_A \in \mathrm{IndHilb}
\rightsquigarrow Q_\bullet^{\mathrm{an}}(A).
\]
Each inclusion is strict: the algebraic VOA is dense in the sewing
envelope, the sewing envelope gives rise to the IndHilb-valued
factorization object, and the analytic shadows $Q_g^{\mathrm{an}}$
are the coderived invariants.  The modular cumulant transform $M(A)$
controls the transition between levels: the generating-function
quintuple $\tau_A$ determines the convergence rate of the sewing
expansion, while the full MC element $\Theta_A$ packages the
algebraic data that the analytic levels complete.
\end{remark}

\subsection{Open directions}
\label{subsec:open-directions-extended}

\begin{question}[Non-standard landscape]
\label{q:non-standard-landscape}
The quantization dictionary is uniform for the standard landscape.
What happens for non-standard families?  Specifically:
\begin{enumerate}[leftmargin=2.4em]
\item Does there exist a filtered quadratic PVA with
$\Mod(A) \neq 0$?

\item Does there exist a PVA with $\dim H^1_{\mathrm{rel}} > 1$?

\item Does there exist a chirally Koszul vertex algebra whose
sewing envelope fails to converge?
\end{enumerate}
Affirmative answers to any of these would bound the reach of the
present programme.
\end{question}

\begin{question}[Lattice sewing beyond even unimodular]
\label{q:lattice-sewing-beyond}
The lattice partition function (Theorem~\ref{thm:lattice-partition-function})
is proved for even unimodular lattices.  For general even lattices
$\Lambda \neq \Lambda^*$, the genus-$g$ theta function
$\Theta_{\Lambda,g}(\Omega)$ must be replaced by a sum over cosets
$\Lambda^*/\Lambda$.  The sewing convergence
(Theorem~\ref{thm:lattice-sewing-convergence}) holds for all even lattices,
but the explicit partition function formula requires the coset
decomposition.
\end{question}

\begin{question}[The non-principal frontier]
\label{q:non-principal-frontier}
The $W_N$ computations above are all for the principal nilpotent
$f_{\mathrm{prin}} \in \mathfrak{sl}_N$.  For non-principal nilpotent
orbits, the $W$-algebra $W(\mathfrak{g}, f)$ may have different
generators and different triangularity properties.  The hook-type
nilpotents in type A provide the first proved non-principal corridor
for Koszul duality
(Vol~I, \S\ref*{sec:nonprincipal-w-duality}).
Extending the modular PVA quantization to non-principal $W$-algebras
requires verifying the triangular vanishing theorem in each case.
\end{question}

\subsection{The Drinfeld--Sokolov route for $W_4$ and $W_5$}
\label{subsec:ds-route-W4-W5}

\begin{theorem}[$W_4$ quantization via Drinfeld--Sokolov reduction;
\ClaimStatusProvedHere]
\label{thm:W4-quantization-ds}
The quantized $W_4$ vertex algebra is obtained by the composite
\[
\text{affine PVA at level }k
\;\xrightarrow{\;\text{Thm~\ref{thm:affine-half-space-bv}}\;}
\;\widehat{\mathfrak{sl}}_4\text{ at level }K_{\mathrm{eff}}
\;\xrightarrow{\;\text{DS reduction}\;}
\;W_4\text{ at }c_{\mathrm{qu}},
\]
where $K_{\mathrm{eff}} = k + \hbar\,\delta_{\mathfrak{sl}_4}$ and
\[
c_{\mathrm{qu}}
=
3\,\frac{(K_{\mathrm{eff}}-1)(4K_{\mathrm{eff}}-1)}{K_{\mathrm{eff}}+4}.
\]
The quantization is one-loop exact: $c_{\mathrm{qu}}$ receives no
corrections beyond the one-loop level shift.
\end{theorem}

\begin{proof}
The affine half-space BV theorem produces the boundary VOA
$V^{K_{\mathrm{eff}}}(\mathfrak{sl}_4)$.  Principal DS reduction
yields $W_4$ at the central charge determined by the Feigin--Frenkel
formula applied to $\mathfrak{sl}_4$.  One-loop exactness is inherited
from the affine case, as in the Virasoro argument
(Theorem~\ref{thm:virasoro-quantization-ds}).
\end{proof}

\begin{theorem}[$W_5$ quantization via Drinfeld--Sokolov reduction;
\ClaimStatusProvedHere]
\label{thm:W5-quantization-ds}
The quantized $W_5$ vertex algebra is obtained by DS reduction from
$\widehat{\mathfrak{sl}}_5$ at the one-loop shifted level, with
\[
c_{\mathrm{qu}}
=
4\,\frac{(K_{\mathrm{eff}}-1)(5K_{\mathrm{eff}}-1)}{K_{\mathrm{eff}}+5}.
\]
\end{theorem}

\begin{proof}
Identical to $W_4$, replacing $\mathfrak{sl}_4$ by $\mathfrak{sl}_5$.
\end{proof}

\begin{proposition}[General $W_N$ central charge formula;
\ClaimStatusProvedHere]
\label{prop:WN-central-charge}
For the principal $W_N$ algebra obtained by DS reduction from
$\widehat{\mathfrak{sl}}_N$ at level $k$:
\begin{equation}\label{eq:WN-central-charge}
c(W_N, k)
=
(N-1)\biggl(1 - \frac{N(N+1)}{(k+N)(k+N-1)}\biggr)
\cdot (k+N-1).
\end{equation}
At the one-loop shifted level $K_{\mathrm{eff}} = k + \hbar\,\delta_N$,
this gives the quantum central charge as a rational function of $k$
and $N$, confirming the one-dimensional modular lift torsor.
\end{proposition}

\begin{proof}
This is the Feigin--Frenkel central charge formula for the principal
$W$-algebra of $\mathfrak{sl}_N$.  It specializes to $c = 1 - 6(k+1)^2/(k+2)$
for $N=2$ (Virasoro) and to the formulas in
Theorems~\ref{thm:W4-quantization-ds} and~\ref{thm:W5-quantization-ds}
for $N=4,5$.
\end{proof}

\begin{remark}[Consistency of the two routes for all $W_N$]
\label{rem:WN-consistency}
For every $N \ge 2$, the direct modular obstruction route
(Proposition~\ref{prop:WN-H1}) and the DS route
(Proposition~\ref{prop:WN-central-charge}) are consistent:
both produce a one-parameter family of quantum $W_N$ vertex algebras
parametrized by a single scalar (the central charge), and the quantum
correction is a shift of that scalar.  The DS route pins down the
exact formula; the obstruction-theoretic route proves uniqueness.
This agreement is a nontrivial check on both the algebraic modular
theory and the DS reduction functor.
\end{remark}

\subsection{Genus-two structure for $W_4$}
\label{subsec:W4-genus-two}

\begin{proposition}[Genus-two forcing for $W_4$; \ClaimStatusProvedHere]
\label{prop:W4-genus-two-forcing}
For the classical $W_4$ family, the reduced genus-two forcing of the
unique relevant genus-one direction $\Xi_c = \partial_c P_c$ is exact:
\[
\mathfrak{F}_2(\Xi_c)
=
\delta_{P_c}(-\partial_c^2 P_c) + \Delta_{\mathrm{cyc}}\partial_c P_c.
\]
Both terms are coboundaries, so the genus-two obstruction vanishes
in the relevant sector.
\end{proposition}

\begin{proof}
The argument is identical to the $W_3$ case
(Corollary~\ref{cor:w3-no-quadratic-self-obstruction}):
$\Delta_{\mathrm{cyc}} P_c = 0$ by triangular vanishing
(Theorem~\ref{thm:W4-mod-vanishing}), and the quadratic
self-interaction is a coboundary by differentiation of
$[P_c, P_c] = 0$.
\end{proof}

\begin{remark}[Genus-two entry point for $W_4$]
The $W_4$ algebra, like $W_3$, exhibits all three genus-two channels
of Remark~\ref{rem:vol2-genus-two-entry}:
$\Theta^{(2)}_{\mathrm{loop}\circ\mathrm{loop}}$ (iterated one-loop),
$\Theta^{(2)}_{\mathrm{sep}\circ\mathrm{loop}}$ (separating degeneration
composed with one-loop), and
$\Theta^{(2)}_{\mathrm{pf}}$ (planted-forest contribution).
The first two channels are forced to vanish in the relevant sector
by the genus-two forcing proposition.  The third channel is the
rigid planted-forest/log-FM sector, which carries the
genuinely new genus-two information.
\end{remark}

\subsection{The affine half-space BV theorem for higher rank}
\label{subsec:affine-higher-rank}

\begin{theorem}[Affine half-space BV at rank $N$; \ClaimStatusConditional]
\label{thm:affine-half-space-bv-rank-N}
Let $\mathfrak{g} = \mathfrak{sl}_N$ and $k$ a non-critical level.
Under Hypotheses~(H1)--(H4) for the $\widehat{\mathfrak{sl}}_N$
Poisson sigma model, the boundary vertex algebra on
$\mathbb{R}_+ \times X$ is $V^{K_{\mathrm{eff}}}(\mathfrak{sl}_N)$,
where
\[
K_{\mathrm{eff}} = k + \hbar\,\delta_N,
\qquad
\delta_N = h^\vee = N.
\]
The one-loop shift is the dual Coxeter number of $\mathfrak{sl}_N$.
\end{theorem}

\begin{remark}[The one-loop shift pattern]
The one-loop shift $\delta_N = h^\vee(\mathfrak{sl}_N) = N$ is a universal
feature of the affine half-space BV construction: it arises from the
ghost determinant in the BV quantization and equals the dual Coxeter
number for any simple Lie algebra.  For the $W_N$ quantization via DS,
this shift propagates to the central charge through the Feigin--Frenkel
formula, producing a quantum correction that is polynomial in $k^{-1}$.
\end{remark}
